%%%%%%%%%%%%%%%%%%%%%%%%%%%%%%%%%%%%%%%%%%%%%%%%%%%%%%%%%%%%
% 8.18 STATISTICAL LEARNING
% The Composition of Advanced Mathematics
%%%%%%%%%%%%%%%%%%%%%%%%%%%%%%%%%%%%%%%%%%%%%%%%%%%%%%%%%%%%

\section{Statistical Learning}
\label{sec:statistical-learning}

\prerequisite{
Probability, statistics, regression, classification, multivariate
statistics, optimization, linear algebra, likelihood, estimation,
cross-validation, and basic computational methods.
}

\learningobjectives{
By the end of this section, the reader should be able to:
\begin{itemize}
    \item distinguish supervised and unsupervised learning;
    \item distinguish regression and classification tasks;
    \item define training, validation, and test data;
    \item understand prediction error and generalization error;
    \item explain the bias--variance tradeoff;
    \item recognize underfitting and overfitting;
    \item understand empirical risk minimization;
    \item understand regularization;
    \item derive ridge-regression estimators;
    \item understand lasso regression;
    \item understand elastic-net regularization;
    \item understand feature selection;
    \item distinguish training error from test error;
    \item understand cross-validation;
    \item understand \(k\)-fold and leave-one-out cross-validation;
    \item understand bootstrap-based model assessment conceptually;
    \item understand classification loss functions;
    \item understand logistic regression as a classification model;
    \item understand linear and quadratic discriminant analysis;
    \item understand \(k\)-nearest-neighbor methods;
    \item understand decision trees;
    \item understand regression trees;
    \item understand classification trees;
    \item understand impurity measures;
    \item understand pruning;
    \item understand bagging;
    \item understand random forests;
    \item understand boosting conceptually;
    \item understand gradient boosting conceptually;
    \item recognize support-vector machines;
    \item understand margins and kernels conceptually;
    \item understand kernel methods;
    \item understand principal components regression;
    \item understand partial least squares conceptually;
    \item recognize clustering methods;
    \item understand \(k\)-means clustering;
    \item understand hierarchical clustering;
    \item understand distance and similarity measures;
    \item understand model complexity and tuning parameters;
    \item understand hyperparameter selection;
    \item understand calibration and discrimination;
    \item understand common performance metrics;
    \item understand imbalanced classification;
    \item understand feature preprocessing and data leakage;
    \item recognize high-dimensional statistical problems;
    \item understand sparsity;
    \item recognize ensemble learning;
    \item understand interpretability and variable importance;
    \item understand fairness and distribution shift conceptually;
    \item connect statistical learning with regression, multivariate
          statistics, optimization, and computational statistics.
\end{itemize}
}

%%%%%%%%%%%%%%%%%%%%%%%%%%%%%%%%%%%%%%%%%%%%%%%%%%%%%%%%%%%%
\subsection{What Is Statistical Learning?}
\label{subsec:what-is-statistical-learning}
%%%%%%%%%%%%%%%%%%%%%%%%%%%%%%%%%%%%%%%%%%%%%%%%%%%%%%%%%%%%

Statistical learning develops methods for discovering patterns,
estimating relationships, and making predictions from data.

A typical dataset consists of predictors

\[
\boxed{
X_1,\ldots,X_p
}
\]

and possibly a response

\[
\boxed{
Y.
}
\]

The goal may be:

\[
\boxed{
\text{prediction},
}
\]

\[
\boxed{
\text{classification},
}
\]

\[
\boxed{
\text{dimension reduction},
}
\]

or

\[
\boxed{
\text{structure discovery}.
}
\]

%%%%%%%%%%%%%%%%%%%%%%%%%%%%%%%%%%%%%%%%%%%%%%%%%%%%%%%%%%%%
\subsection{Supervised Learning}
\label{subsec:supervised-learning}
%%%%%%%%%%%%%%%%%%%%%%%%%%%%%%%%%%%%%%%%%%%%%%%%%%%%%%%%%%%%

In supervised learning, each training observation includes predictors and
a response:

\[
\boxed{
(\mathbf{x}_i,y_i).
}
\]

The objective is to learn a function

\[
\boxed{
\widehat{f}(\mathbf{x})
}
\]

that predicts \(Y\) from \(\mathbf{X}\).

%%%%%%%%%%%%%%%%%%%%%%%%%%%%%%%%%%%%%%%%%%%%%%%%%%%%%%%%%%%%
\subsection{Unsupervised Learning}
\label{subsec:unsupervised-learning}
%%%%%%%%%%%%%%%%%%%%%%%%%%%%%%%%%%%%%%%%%%%%%%%%%%%%%%%%%%%%

In unsupervised learning, no response variable is specified.

The observations are

\[
\boxed{
\mathbf{x}_1,\ldots,\mathbf{x}_n.
}
\]

The goal is to discover structure such as:

\[
\boxed{
\text{clusters},
}
\]

\[
\boxed{
\text{latent directions},
}
\]

or

\[
\boxed{
\text{low-dimensional representations}.
}
\]

%%%%%%%%%%%%%%%%%%%%%%%%%%%%%%%%%%%%%%%%%%%%%%%%%%%%%%%%%%%%
\subsection{Regression Problems}
\label{subsec:statistical-learning-regression}
%%%%%%%%%%%%%%%%%%%%%%%%%%%%%%%%%%%%%%%%%%%%%%%%%%%%%%%%%%%%

Regression predicts a quantitative response.

Examples include:

\[
\boxed{
\text{income},
}
\]

\[
\boxed{
\text{temperature},
}
\]

\[
\boxed{
\text{sales},
}
\]

or

\[
\boxed{
\text{disease duration}.
}
\]

A regression model aims to estimate

\[
\boxed{
f(\mathbf{x})
=
\mathbb{E}
[
Y\mid\mathbf{X}=\mathbf{x}
].
}
\]

%%%%%%%%%%%%%%%%%%%%%%%%%%%%%%%%%%%%%%%%%%%%%%%%%%%%%%%%%%%%
\subsection{Classification Problems}
\label{subsec:statistical-learning-classification}
%%%%%%%%%%%%%%%%%%%%%%%%%%%%%%%%%%%%%%%%%%%%%%%%%%%%%%%%%%%%

Classification predicts a categorical response.

For binary classification,

\[
\boxed{
Y\in\{0,1\}.
}
\]

A probabilistic classifier may estimate

\[
\boxed{
P
(
Y=1
\mid
\mathbf{X}=\mathbf{x}
).
}
\]

A class label is then chosen using a decision rule.

%%%%%%%%%%%%%%%%%%%%%%%%%%%%%%%%%%%%%%%%%%%%%%%%%%%%%%%%%%%%
\subsection{Prediction Versus Explanation}
\label{subsec:prediction-versus-explanation}
%%%%%%%%%%%%%%%%%%%%%%%%%%%%%%%%%%%%%%%%%%%%%%%%%%%%%%%%%%%%

A model that predicts well does not necessarily provide a correct causal
or scientific explanation.

Prediction focuses on future accuracy.

Scientific inference may instead focus on:

\[
\boxed{
\text{effect estimation},
}
\]

\[
\boxed{
\text{uncertainty},
}
\]

or

\[
\boxed{
\text{causal interpretation}.
}
\]

These goals can require different methods.

%%%%%%%%%%%%%%%%%%%%%%%%%%%%%%%%%%%%%%%%%%%%%%%%%%%%%%%%%%%%
\subsection{Training Data}
\label{subsec:training-data}
%%%%%%%%%%%%%%%%%%%%%%%%%%%%%%%%%%%%%%%%%%%%%%%%%%%%%%%%%%%%

The training set is used to estimate model parameters.

Write

\[
\boxed{
\mathcal{D}_{train}.
}
\]

For supervised learning,

\[
\mathcal{D}_{train}
=
\{
(\mathbf{x}_i,y_i)
\}_{i=1}^{n}.
\]

The fitted model is determined from these observations.

%%%%%%%%%%%%%%%%%%%%%%%%%%%%%%%%%%%%%%%%%%%%%%%%%%%%%%%%%%%%
\subsection{Validation Data}
\label{subsec:validation-data}
%%%%%%%%%%%%%%%%%%%%%%%%%%%%%%%%%%%%%%%%%%%%%%%%%%%%%%%%%%%%

A validation set is used to compare model settings or tuning parameters.

For example, validation data may determine:

\[
\boxed{
\lambda
}
\]

in ridge regression or

\[
\boxed{
k
}
\]

in \(k\)-nearest-neighbor classification.

Validation data should not be repeatedly reused as if they were unseen
test data.

%%%%%%%%%%%%%%%%%%%%%%%%%%%%%%%%%%%%%%%%%%%%%%%%%%%%%%%%%%%%
\subsection{Test Data}
\label{subsec:test-data}
%%%%%%%%%%%%%%%%%%%%%%%%%%%%%%%%%%%%%%%%%%%%%%%%%%%%%%%%%%%%

The test set is reserved for final evaluation.

It should represent observations not used to:

\[
\boxed{
\text{fit the model},
}
\]

\[
\boxed{
\text{choose predictors},
}
\]

or

\[
\boxed{
\text{tune hyperparameters}.
}
\]

The test error estimates generalization performance.

%%%%%%%%%%%%%%%%%%%%%%%%%%%%%%%%%%%%%%%%%%%%%%%%%%%%%%%%%%%%
\subsection{Generalization}
\label{subsec:generalization}
%%%%%%%%%%%%%%%%%%%%%%%%%%%%%%%%%%%%%%%%%%%%%%%%%%%%%%%%%%%%

Generalization refers to how well a learned model performs on new data
from the target population.

A model with excellent training performance may generalize poorly.

Thus the key quantity is not merely

\[
\boxed{
\text{training error},
}
\]

but

\[
\boxed{
\text{future prediction error}.
}
\]

%%%%%%%%%%%%%%%%%%%%%%%%%%%%%%%%%%%%%%%%%%%%%%%%%%%%%%%%%%%%
\subsection{Loss Functions}
\label{subsec:learning-loss-functions}
%%%%%%%%%%%%%%%%%%%%%%%%%%%%%%%%%%%%%%%%%%%%%%%%%%%%%%%%%%%%

A loss function measures the cost of a prediction.

For regression, squared-error loss is

\[
\boxed{
L
(
y,\widehat{y}
)
=
(
y-\widehat{y}
)^2.
}
\]

Absolute-error loss is

\[
\boxed{
L
(
y,\widehat{y}
)
=
|
y-\widehat{y}
|.
}
\]

%%%%%%%%%%%%%%%%%%%%%%%%%%%%%%%%%%%%%%%%%%%%%%%%%%%%%%%%%%%%
\subsection{Classification Loss}
\label{subsec:classification-loss}
%%%%%%%%%%%%%%%%%%%%%%%%%%%%%%%%%%%%%%%%%%%%%%%%%%%%%%%%%%%%

For classification, zero-one loss is

\[
\boxed{
L
(
y,\widehat{y}
)
=
\mathbf{1}_{\{y\neq\widehat{y}\}}.
}
\]

The average zero-one loss is the misclassification rate.

Other losses may reflect unequal costs of different mistakes.

%%%%%%%%%%%%%%%%%%%%%%%%%%%%%%%%%%%%%%%%%%%%%%%%%%%%%%%%%%%%
\subsection{Risk}
\label{subsec:learning-risk}
%%%%%%%%%%%%%%%%%%%%%%%%%%%%%%%%%%%%%%%%%%%%%%%%%%%%%%%%%%%%

The population prediction risk is

\[
\boxed{
R(f)
=
\mathbb{E}
[
L
(
Y,f(\mathbf{X})
)
].
}
\]

The expectation is taken over future predictor-response pairs.

The ideal predictor minimizes population risk.

%%%%%%%%%%%%%%%%%%%%%%%%%%%%%%%%%%%%%%%%%%%%%%%%%%%%%%%%%%%%
\subsection{Empirical Risk}
\label{subsec:empirical-risk}
%%%%%%%%%%%%%%%%%%%%%%%%%%%%%%%%%%%%%%%%%%%%%%%%%%%%%%%%%%%%

The empirical risk is the average training loss:

\[
\boxed{
\widehat{R}_n(f)
=
\frac1n
\sum_{i=1}^{n}
L
(
Y_i,
f(\mathbf{X}_i)
).
}
\]

Many learning algorithms choose a model that minimizes empirical risk or
a penalized version of it.

%%%%%%%%%%%%%%%%%%%%%%%%%%%%%%%%%%%%%%%%%%%%%%%%%%%%%%%%%%%%
\subsection{Empirical Risk Minimization}
\label{subsec:empirical-risk-minimization}
%%%%%%%%%%%%%%%%%%%%%%%%%%%%%%%%%%%%%%%%%%%%%%%%%%%%%%%%%%%%

Given a model class

\[
\mathcal{F},
\]

empirical risk minimization chooses

\[
\boxed{
\widehat{f}
=
\arg\min_{f\in\mathcal{F}}
\widehat{R}_n(f).
}
\]

A highly flexible class can reduce training error dramatically.

However, minimizing training error alone can cause overfitting.

%%%%%%%%%%%%%%%%%%%%%%%%%%%%%%%%%%%%%%%%%%%%%%%%%%%%%%%%%%%%
\subsection{Irreducible Error}
\label{subsec:irreducible-error}
%%%%%%%%%%%%%%%%%%%%%%%%%%%%%%%%%%%%%%%%%%%%%%%%%%%%%%%%%%%%

Suppose

\[
\boxed{
Y
=
f(\mathbf{X})
+
\varepsilon,
}
\]

with

\[
\mathbb{E}
[
\varepsilon\mid\mathbf{X}
]
=
0.
\]

Even if \(f\) were known exactly, future observations still contain
random noise.

This is irreducible error.

%%%%%%%%%%%%%%%%%%%%%%%%%%%%%%%%%%%%%%%%%%%%%%%%%%%%%%%%%%%%
\subsection{Model Error}
\label{subsec:model-error}
%%%%%%%%%%%%%%%%%%%%%%%%%%%%%%%%%%%%%%%%%%%%%%%%%%%%%%%%%%%%

Prediction error above irreducible noise may arise from:

\[
\boxed{
\text{bias},
}
\]

\[
\boxed{
\text{variance},
}
\]

or

\[
\boxed{
\text{model misspecification}.
}
\]

The goal of statistical learning is to control these sources of error.

%%%%%%%%%%%%%%%%%%%%%%%%%%%%%%%%%%%%%%%%%%%%%%%%%%%%%%%%%%%%
\subsection{Bias--Variance Tradeoff}
\label{subsec:bias-variance-tradeoff-learning}
%%%%%%%%%%%%%%%%%%%%%%%%%%%%%%%%%%%%%%%%%%%%%%%%%%%%%%%%%%%%

For squared-error prediction at a fixed point \(\mathbf{x}_0\), a
classical decomposition is

\[
\boxed{
\mathbb{E}
\left[
(
Y_0
-
\widehat{f}(\mathbf{x}_0)
)^2
\right]
=
\sigma^2
+
\operatorname{Bias}
[
\widehat{f}(\mathbf{x}_0)
]^2
+
\operatorname{Var}
[
\widehat{f}(\mathbf{x}_0)
].
}
\]

The first term is irreducible noise.

%%%%%%%%%%%%%%%%%%%%%%%%%%%%%%%%%%%%%%%%%%%%%%%%%%%%%%%%%%%%
\subsection{Flexible Versus Rigid Models}
\label{subsec:flexible-versus-rigid}
%%%%%%%%%%%%%%%%%%%%%%%%%%%%%%%%%%%%%%%%%%%%%%%%%%%%%%%%%%%%

A highly rigid model may have:

\[
\boxed{
\text{high bias}
}
\]

and

\[
\boxed{
\text{low variance}.
}
\]

A highly flexible model may have:

\[
\boxed{
\text{low bias}
}
\]

and

\[
\boxed{
\text{high variance}.
}
\]

Good predictive performance requires balancing these effects.

%%%%%%%%%%%%%%%%%%%%%%%%%%%%%%%%%%%%%%%%%%%%%%%%%%%%%%%%%%%%
\subsection{Underfitting}
\label{subsec:underfitting}
%%%%%%%%%%%%%%%%%%%%%%%%%%%%%%%%%%%%%%%%%%%%%%%%%%%%%%%%%%%%

Underfitting occurs when the model is too simple to capture important
structure.

Typical symptoms include:

\[
\boxed{
\text{high training error}
}
\]

and

\[
\boxed{
\text{high test error}.
}
\]

Adding relevant flexibility may improve performance.

%%%%%%%%%%%%%%%%%%%%%%%%%%%%%%%%%%%%%%%%%%%%%%%%%%%%%%%%%%%%
\subsection{Overfitting}
\label{subsec:overfitting}
%%%%%%%%%%%%%%%%%%%%%%%%%%%%%%%%%%%%%%%%%%%%%%%%%%%%%%%%%%%%

Overfitting occurs when a model adapts too strongly to random noise or
peculiarities of the training data.

Typical symptoms include:

\[
\boxed{
\text{very low training error}
}
\]

but

\[
\boxed{
\text{high test error}.
}
\]

Regularization and validation are major defenses against overfitting.

%%%%%%%%%%%%%%%%%%%%%%%%%%%%%%%%%%%%%%%%%%%%%%%%%%%%%%%%%%%%
\subsection{Model Complexity}
\label{subsec:model-complexity-learning}
%%%%%%%%%%%%%%%%%%%%%%%%%%%%%%%%%%%%%%%%%%%%%%%%%%%%%%%%%%%%

Model complexity can be controlled through:

\[
\boxed{
\text{number of predictors},
}
\]

\[
\boxed{
\text{tree depth},
}
\]

\[
\boxed{
\text{polynomial degree},
}
\]

\[
\boxed{
\text{regularization strength},
}
\]

or other hyperparameters.

The best complexity is generally determined using out-of-sample
performance.

%%%%%%%%%%%%%%%%%%%%%%%%%%%%%%%%%%%%%%%%%%%%%%%%%%%%%%%%%%%%
\subsection{Training and Test Error Curves}
\label{subsec:training-test-error-curves}
%%%%%%%%%%%%%%%%%%%%%%%%%%%%%%%%%%%%%%%%%%%%%%%%%%%%%%%%%%%%

As flexibility increases, training error generally decreases.

Test error often follows a U-shaped pattern:

\[
\boxed{
\text{underfitting}
\rightarrow
\text{optimal complexity}
\rightarrow
\text{overfitting}.
}
\]

The minimum test or validation error determines a useful model complexity
for prediction.

%%%%%%%%%%%%%%%%%%%%%%%%%%%%%%%%%%%%%%%%%%%%%%%%%%%%%%%%%%%%
\subsection{Regularization}
\label{subsec:regularization}
%%%%%%%%%%%%%%%%%%%%%%%%%%%%%%%%%%%%%%%%%%%%%%%%%%%%%%%%%%%%

Regularization modifies the fitting objective by penalizing model
complexity.

A generic penalized objective is

\[
\boxed{
\text{Loss}
+
\lambda
\cdot
\text{Penalty}.
}
\]

The tuning parameter

\[
\boxed{
\lambda\geq0
}
\]

controls the strength of regularization.

%%%%%%%%%%%%%%%%%%%%%%%%%%%%%%%%%%%%%%%%%%%%%%%%%%%%%%%%%%%%
\subsection{Ridge Regression}
\label{subsec:ridge-regression}
%%%%%%%%%%%%%%%%%%%%%%%%%%%%%%%%%%%%%%%%%%%%%%%%%%%%%%%%%%%%

Ridge regression minimizes

\[
\boxed{
\sum_{i=1}^{n}
(
y_i
-
\beta_0
-
\mathbf{x}_i^T
\boldsymbol{\beta}
)^2
+
\lambda
\sum_{j=1}^{p}
\beta_j^2.
}
\]

The intercept is usually not penalized.

The penalty shrinks coefficients toward zero.

%%%%%%%%%%%%%%%%%%%%%%%%%%%%%%%%%%%%%%%%%%%%%%%%%%%%%%%%%%%%
\subsection{Matrix Form of Ridge Regression}
\label{subsec:ridge-matrix-form}
%%%%%%%%%%%%%%%%%%%%%%%%%%%%%%%%%%%%%%%%%%%%%%%%%%%%%%%%%%%%

After centering the response and predictors so an intercept is not
required, ridge minimizes

\[
\boxed{
\|
\mathbf{y}
-
X\boldsymbol{\beta}
\|_2^2
+
\lambda
\|
\boldsymbol{\beta}
\|_2^2.
}
\]

Differentiate with respect to \(\boldsymbol{\beta}\):

\[
-2X^T
(
\mathbf{y}
-
X\boldsymbol{\beta}
)
+
2\lambda
\boldsymbol{\beta}
=
0.
\]

Therefore,

\[
\boxed{
(
X^TX+\lambda I
)
\widehat{\boldsymbol{\beta}}_{ridge}
=
X^T\mathbf{y}.
}
\]

Hence,

\[
\boxed{
\widehat{\boldsymbol{\beta}}_{ridge}
=
(
X^TX+\lambda I
)^{-1}
X^T\mathbf{y}.
}
\]

%%%%%%%%%%%%%%%%%%%%%%%%%%%%%%%%%%%%%%%%%%%%%%%%%%%%%%%%%%%%
\subsection{Why Ridge Helps}
\label{subsec:why-ridge-helps}
%%%%%%%%%%%%%%%%%%%%%%%%%%%%%%%%%%%%%%%%%%%%%%%%%%%%%%%%%%%%

Ordinary least squares can have high variance when predictors are highly
correlated or when

\[
p
\]

is large relative to \(n\).

Adding

\[
\lambda I
\]

stabilizes the matrix inversion.

Ridge accepts some bias in exchange for lower variance.

%%%%%%%%%%%%%%%%%%%%%%%%%%%%%%%%%%%%%%%%%%%%%%%%%%%%%%%%%%%%
\subsection{Ridge and Multicollinearity}
\label{subsec:ridge-multicollinearity}
%%%%%%%%%%%%%%%%%%%%%%%%%%%%%%%%%%%%%%%%%%%%%%%%%%%%%%%%%%%%

If

\[
X^TX
\]

has very small eigenvalues, ordinary least-squares coefficients may be
unstable.

Ridge replaces an eigenvalue

\[
d_j
\]

conceptually with

\[
\boxed{
d_j+\lambda.
}
\]

This reduces amplification along poorly identified directions.

%%%%%%%%%%%%%%%%%%%%%%%%%%%%%%%%%%%%%%%%%%%%%%%%%%%%%%%%%%%%
\subsection{Worked Example: Ridge Shrinkage}
\label{subsec:worked-ridge-shrinkage}
%%%%%%%%%%%%%%%%%%%%%%%%%%%%%%%%%%%%%%%%%%%%%%%%%%%%%%%%%%%%

\begin{workedexamplebox}{Effect of the Penalty}

Suppose a centered one-predictor problem has

\[
X^TX=10
\]

and

\[
X^Ty=20.
\]

The ordinary least-squares coefficient is

\[
\widehat{\beta}_{OLS}
=
\frac{20}{10}
=
2.
\]

If

\[
\lambda=5,
\]

then

\[
\widehat{\beta}_{ridge}
=
\frac{
20
}{
10+5
}.
\]

Therefore,

\begin{answerbox}
\[
\boxed{
\widehat{\beta}_{ridge}
=
\frac43
\approx
1.333.
}
\]

The coefficient is shrunk toward zero.
\end{answerbox}

\end{workedexamplebox}

%%%%%%%%%%%%%%%%%%%%%%%%%%%%%%%%%%%%%%%%%%%%%%%%%%%%%%%%%%%%
\subsection{Standardization Before Regularization}
\label{subsec:standardization-regularization}
%%%%%%%%%%%%%%%%%%%%%%%%%%%%%%%%%%%%%%%%%%%%%%%%%%%%%%%%%%%%

Regularization depends on coefficient magnitude.

If predictors have very different units, the penalty can affect them
unequally.

Therefore predictors are often standardized:

\[
\boxed{
Z_j
=
\frac{
X_j-\overline{X}_j
}{
s_j
}.
}
\]

This places predictors on comparable scales before penalization.

%%%%%%%%%%%%%%%%%%%%%%%%%%%%%%%%%%%%%%%%%%%%%%%%%%%%%%%%%%%%
\subsection{Lasso Regression}
\label{subsec:lasso-regression}
%%%%%%%%%%%%%%%%%%%%%%%%%%%%%%%%%%%%%%%%%%%%%%%%%%%%%%%%%%%%

The lasso minimizes

\[
\boxed{
\sum_{i=1}^{n}
(
y_i
-
\beta_0
-
\mathbf{x}_i^T
\boldsymbol{\beta}
)^2
+
\lambda
\sum_{j=1}^{p}
|\beta_j|.
}
\]

The penalty is the

\[
\ell_1
\]

norm of the coefficient vector.

%%%%%%%%%%%%%%%%%%%%%%%%%%%%%%%%%%%%%%%%%%%%%%%%%%%%%%%%%%%%
\subsection{Lasso Sparsity}
\label{subsec:lasso-sparsity}
%%%%%%%%%%%%%%%%%%%%%%%%%%%%%%%%%%%%%%%%%%%%%%%%%%%%%%%%%%%%

Unlike ridge regression, the lasso can set coefficients exactly equal to
zero:

\[
\boxed{
\widehat{\beta}_j=0.
}
\]

Thus the lasso performs both:

\[
\boxed{
\text{shrinkage}
}
\]

and

\[
\boxed{
\text{variable selection}.
}
\]

%%%%%%%%%%%%%%%%%%%%%%%%%%%%%%%%%%%%%%%%%%%%%%%%%%%%%%%%%%%%
\subsection{Ridge Versus Lasso}
\label{subsec:ridge-versus-lasso}
%%%%%%%%%%%%%%%%%%%%%%%%%%%%%%%%%%%%%%%%%%%%%%%%%%%%%%%%%%%%

Ridge uses

\[
\boxed{
\sum_j\beta_j^2.
}
\]

Lasso uses

\[
\boxed{
\sum_j|\beta_j|.
}
\]

Ridge usually retains all predictors with reduced coefficients.

Lasso can produce sparse models with some coefficients exactly zero.

%%%%%%%%%%%%%%%%%%%%%%%%%%%%%%%%%%%%%%%%%%%%%%%%%%%%%%%%%%%%
\subsection{Constraint Form of Ridge}
\label{subsec:ridge-constraint-form}
%%%%%%%%%%%%%%%%%%%%%%%%%%%%%%%%%%%%%%%%%%%%%%%%%%%%%%%%%%%%

Ridge can also be written as

\[
\boxed{
\min_{\boldsymbol{\beta}}
\|
\mathbf{y}
-
X\boldsymbol{\beta}
\|_2^2
}
\]

subject to

\[
\boxed{
\sum_j
\beta_j^2
\leq
t.
}
\]

The tuning parameter \(t\) corresponds to a penalty value \(\lambda\).

%%%%%%%%%%%%%%%%%%%%%%%%%%%%%%%%%%%%%%%%%%%%%%%%%%%%%%%%%%%%
\subsection{Constraint Form of Lasso}
\label{subsec:lasso-constraint-form}
%%%%%%%%%%%%%%%%%%%%%%%%%%%%%%%%%%%%%%%%%%%%%%%%%%%%%%%%%%%%

Similarly, lasso can be written

\[
\boxed{
\min_{\boldsymbol{\beta}}
\|
\mathbf{y}
-
X\boldsymbol{\beta}
\|_2^2
}
\]

subject to

\[
\boxed{
\sum_j
|\beta_j|
\leq
t.
}
\]

The geometry of the \(\ell_1\) constraint helps explain why sparse
solutions arise.

%%%%%%%%%%%%%%%%%%%%%%%%%%%%%%%%%%%%%%%%%%%%%%%%%%%%%%%%%%%%
\subsection{Elastic Net}
\label{subsec:elastic-net}
%%%%%%%%%%%%%%%%%%%%%%%%%%%%%%%%%%%%%%%%%%%%%%%%%%%%%%%%%%%%

Elastic net combines ridge and lasso penalties:

\[
\boxed{
\text{Loss}
+
\lambda
\left[
\alpha
\sum_j|\beta_j|
+
(1-\alpha)
\sum_j\beta_j^2
\right].
}
\]

Here,

\[
\boxed{
0\leq\alpha\leq1.
}
\]

When

\[
\alpha=1,
\]

the penalty is lasso-like.

When

\[
\alpha=0,
\]

it is ridge-like.

%%%%%%%%%%%%%%%%%%%%%%%%%%%%%%%%%%%%%%%%%%%%%%%%%%%%%%%%%%%%
\subsection{Why Elastic Net?}
\label{subsec:why-elastic-net}
%%%%%%%%%%%%%%%%%%%%%%%%%%%%%%%%%%%%%%%%%%%%%%%%%%%%%%%%%%%%

Elastic net can be useful when:

\[
\boxed{
\text{many predictors are correlated},
}
\]

and

\[
\boxed{
\text{sparse solutions are still desired}.
}
\]

It combines grouping behavior from ridge with sparsity from lasso.

%%%%%%%%%%%%%%%%%%%%%%%%%%%%%%%%%%%%%%%%%%%%%%%%%%%%%%%%%%%%
\subsection{Feature Selection}
\label{subsec:feature-selection}
%%%%%%%%%%%%%%%%%%%%%%%%%%%%%%%%%%%%%%%%%%%%%%%%%%%%%%%%%%%%

Feature selection chooses a subset of predictors.

Possible goals include:

\[
\boxed{
\text{better prediction},
}
\]

\[
\boxed{
\text{reduced variance},
}
\]

\[
\boxed{
\text{lower computational cost},
}
\]

or

\[
\boxed{
\text{improved interpretability}.
}
\]

Selection itself creates additional uncertainty.

%%%%%%%%%%%%%%%%%%%%%%%%%%%%%%%%%%%%%%%%%%%%%%%%%%%%%%%%%%%%
\subsection{Best Subset Selection}
\label{subsec:best-subset-selection}
%%%%%%%%%%%%%%%%%%%%%%%%%%%%%%%%%%%%%%%%%%%%%%%%%%%%%%%%%%%%

Best subset selection compares predictor subsets and chooses the best
according to a criterion.

With \(p\) predictors, there are

\[
\boxed{
2^p
}
\]

possible subsets.

Thus exhaustive search becomes computationally difficult as \(p\)
grows.

%%%%%%%%%%%%%%%%%%%%%%%%%%%%%%%%%%%%%%%%%%%%%%%%%%%%%%%%%%%%
\subsection{Forward Selection}
\label{subsec:forward-selection}
%%%%%%%%%%%%%%%%%%%%%%%%%%%%%%%%%%%%%%%%%%%%%%%%%%%%%%%%%%%%

Forward selection begins with a small model and repeatedly adds the
predictor producing the greatest improvement according to a chosen
criterion.

It is computationally cheaper than exhaustive subset search.

However, an early poor choice may not be corrected later.

%%%%%%%%%%%%%%%%%%%%%%%%%%%%%%%%%%%%%%%%%%%%%%%%%%%%%%%%%%%%
\subsection{Backward Elimination}
\label{subsec:backward-elimination}
%%%%%%%%%%%%%%%%%%%%%%%%%%%%%%%%%%%%%%%%%%%%%%%%%%%%%%%%%%%%

Backward elimination begins with a large model and repeatedly removes
predictors.

It requires that the initial model be estimable.

Thus it can be problematic when

\[
p\geq n.
\]

%%%%%%%%%%%%%%%%%%%%%%%%%%%%%%%%%%%%%%%%%%%%%%%%%%%%%%%%%%%%
\subsection{Cross-Validation}
\label{subsec:cross-validation}
%%%%%%%%%%%%%%%%%%%%%%%%%%%%%%%%%%%%%%%%%%%%%%%%%%%%%%%%%%%%

Cross-validation estimates prediction performance by repeatedly fitting a
model on part of the data and evaluating it on held-out observations.

This approximates the error the model may experience on future data.

%%%%%%%%%%%%%%%%%%%%%%%%%%%%%%%%%%%%%%%%%%%%%%%%%%%%%%%%%%%%
\subsection{\(k\)-Fold Cross-Validation}
\label{subsec:k-fold-cross-validation}
%%%%%%%%%%%%%%%%%%%%%%%%%%%%%%%%%%%%%%%%%%%%%%%%%%%%%%%%%%%%

Partition the data into \(k\) folds:

\[
\boxed{
F_1,\ldots,F_k.
}
\]

For each fold:

\begin{enumerate}
    \item fit the model using the other \(k-1\) folds;
    \item predict the held-out fold;
    \item calculate loss.
\end{enumerate}

The cross-validation estimate is the average held-out loss.

%%%%%%%%%%%%%%%%%%%%%%%%%%%%%%%%%%%%%%%%%%%%%%%%%%%%%%%%%%%%
\subsection{\(k\)-Fold CV Estimate}
\label{subsec:kfold-cv-estimate}
%%%%%%%%%%%%%%%%%%%%%%%%%%%%%%%%%%%%%%%%%%%%%%%%%%%%%%%%%%%%

If

\[
\widehat{f}^{(-j)}
\]

denotes the model fitted without fold \(j\), then

\[
\boxed{
CV_k
=
\frac1n
\sum_{j=1}^{k}
\sum_{i\in F_j}
L
\left(
y_i,
\widehat{f}^{(-j)}(\mathbf{x}_i)
\right).
}
\]

This can be calculated for different model complexities.

%%%%%%%%%%%%%%%%%%%%%%%%%%%%%%%%%%%%%%%%%%%%%%%%%%%%%%%%%%%%
\subsection{Leave-One-Out Cross-Validation}
\label{subsec:loocv}
%%%%%%%%%%%%%%%%%%%%%%%%%%%%%%%%%%%%%%%%%%%%%%%%%%%%%%%%%%%%

Leave-one-out cross-validation, abbreviated LOOCV, uses

\[
\boxed{
k=n.
}
\]

Each observation is held out once.

The model is fitted on the remaining

\[
n-1
\]

observations.

LOOCV has low training-set bias but may have high variance as an error
estimator and can be computationally expensive.

%%%%%%%%%%%%%%%%%%%%%%%%%%%%%%%%%%%%%%%%%%%%%%%%%%%%%%%%%%%%
\subsection{Cross-Validation for Tuning}
\label{subsec:cv-for-tuning}
%%%%%%%%%%%%%%%%%%%%%%%%%%%%%%%%%%%%%%%%%%%%%%%%%%%%%%%%%%%%

Suppose a model depends on hyperparameter

\[
\lambda.
\]

Compute

\[
CV(\lambda)
\]

for several values.

Then choose

\[
\boxed{
\widehat{\lambda}
=
\arg\min_\lambda
CV(\lambda).
}
\]

The model is then refit using the chosen value.

%%%%%%%%%%%%%%%%%%%%%%%%%%%%%%%%%%%%%%%%%%%%%%%%%%%%%%%%%%%%
\subsection{Nested Cross-Validation}
\label{subsec:nested-cross-validation}
%%%%%%%%%%%%%%%%%%%%%%%%%%%%%%%%%%%%%%%%%%%%%%%%%%%%%%%%%%%%

If cross-validation is used both for tuning and reporting final
performance, optimistic bias can arise.

Nested cross-validation uses:

\[
\boxed{
\text{inner folds for tuning}
}
\]

and

\[
\boxed{
\text{outer folds for evaluation}.
}
\]

This better separates model selection from model assessment.

%%%%%%%%%%%%%%%%%%%%%%%%%%%%%%%%%%%%%%%%%%%%%%%%%%%%%%%%%%%%
\subsection{Data Leakage}
\label{subsec:data-leakage}
%%%%%%%%%%%%%%%%%%%%%%%%%%%%%%%%%%%%%%%%%%%%%%%%%%%%%%%%%%%%

Data leakage occurs when information from validation or test observations
influences model training.

Examples include:

\[
\boxed{
\text{standardizing using all observations},
}
\]

\[
\boxed{
\text{feature selection before splitting},
}
\]

or

\[
\boxed{
\text{imputing using test-set statistics}.
}
\]

Leakage produces overly optimistic estimates of performance.

%%%%%%%%%%%%%%%%%%%%%%%%%%%%%%%%%%%%%%%%%%%%%%%%%%%%%%%%%%%%
\subsection{Preprocessing Within Resampling}
\label{subsec:preprocessing-within-resampling}
%%%%%%%%%%%%%%%%%%%%%%%%%%%%%%%%%%%%%%%%%%%%%%%%%%%%%%%%%%%%

Any data-dependent preprocessing should be estimated using only the
training portion of each resampling iteration.

This includes:

\[
\boxed{
\text{scaling},
}
\]

\[
\boxed{
\text{imputation},
}
\]

\[
\boxed{
\text{feature selection},
}
\]

and

\[
\boxed{
\text{dimension reduction}.
}
\]

%%%%%%%%%%%%%%%%%%%%%%%%%%%%%%%%%%%%%%%%%%%%%%%%%%%%%%%%%%%%
\subsection{Classification Accuracy}
\label{subsec:classification-accuracy}
%%%%%%%%%%%%%%%%%%%%%%%%%%%%%%%%%%%%%%%%%%%%%%%%%%%%%%%%%%%%

For predicted classes

\[
\widehat{Y}_i,
\]

accuracy is

\[
\boxed{
\text{Accuracy}
=
\frac{
\#\{\widehat{Y}_i=Y_i\}
}{
n
}.
}
\]

Misclassification rate is

\[
\boxed{
1-\text{Accuracy}.
}
\]

Accuracy can be misleading when class frequencies are highly imbalanced.

%%%%%%%%%%%%%%%%%%%%%%%%%%%%%%%%%%%%%%%%%%%%%%%%%%%%%%%%%%%%
\subsection{Confusion Matrix}
\label{subsec:confusion-matrix-learning}
%%%%%%%%%%%%%%%%%%%%%%%%%%%%%%%%%%%%%%%%%%%%%%%%%%%%%%%%%%%%

For binary classification:

\[
\begin{array}{c|cc}
&
\text{Predicted }1
&
\text{Predicted }0
\\
\hline
\text{Actual }1
&
TP
&
FN
\\
\text{Actual }0
&
FP
&
TN
\end{array}
\]

where:

\[
TP
=
\text{true positives},
\]

\[
TN
=
\text{true negatives},
\]

\[
FP
=
\text{false positives},
\]

and

\[
FN
=
\text{false negatives}.
\]

%%%%%%%%%%%%%%%%%%%%%%%%%%%%%%%%%%%%%%%%%%%%%%%%%%%%%%%%%%%%
\subsection{Sensitivity}
\label{subsec:sensitivity-learning}
%%%%%%%%%%%%%%%%%%%%%%%%%%%%%%%%%%%%%%%%%%%%%%%%%%%%%%%%%%%%

Sensitivity, also called recall or true-positive rate, is

\[
\boxed{
\text{Sensitivity}
=
\frac{
TP
}{
TP+FN
}.
}
\]

It measures the proportion of actual positive observations correctly
identified.

%%%%%%%%%%%%%%%%%%%%%%%%%%%%%%%%%%%%%%%%%%%%%%%%%%%%%%%%%%%%
\subsection{Specificity}
\label{subsec:specificity-learning}
%%%%%%%%%%%%%%%%%%%%%%%%%%%%%%%%%%%%%%%%%%%%%%%%%%%%%%%%%%%%

Specificity is

\[
\boxed{
\text{Specificity}
=
\frac{
TN
}{
TN+FP
}.
}
\]

It measures the proportion of actual negative observations correctly
identified.

%%%%%%%%%%%%%%%%%%%%%%%%%%%%%%%%%%%%%%%%%%%%%%%%%%%%%%%%%%%%
\subsection{Precision}
\label{subsec:precision-learning}
%%%%%%%%%%%%%%%%%%%%%%%%%%%%%%%%%%%%%%%%%%%%%%%%%%%%%%%%%%%%

Precision, also called positive predictive value, is

\[
\boxed{
\text{Precision}
=
\frac{
TP
}{
TP+FP
}.
}
\]

It measures the proportion of predicted positives that are actually
positive.

%%%%%%%%%%%%%%%%%%%%%%%%%%%%%%%%%%%%%%%%%%%%%%%%%%%%%%%%%%%%
\subsection{\(F_1\) Score}
\label{subsec:f1-score}
%%%%%%%%%%%%%%%%%%%%%%%%%%%%%%%%%%%%%%%%%%%%%%%%%%%%%%%%%%%%

The \(F_1\) score is the harmonic mean of precision and recall:

\[
\boxed{
F_1
=
2
\frac{
\text{Precision}
\cdot
\text{Recall}
}{
\text{Precision}
+
\text{Recall}
}.
}
\]

It is useful when both false positives and false negatives matter.

%%%%%%%%%%%%%%%%%%%%%%%%%%%%%%%%%%%%%%%%%%%%%%%%%%%%%%%%%%%%
\subsection{Classification Threshold}
\label{subsec:classification-threshold}
%%%%%%%%%%%%%%%%%%%%%%%%%%%%%%%%%%%%%%%%%%%%%%%%%%%%%%%%%%%%

Suppose a model estimates

\[
\widehat{p}(\mathbf{x})
=
P
(
Y=1
\mid
\mathbf{X}=\mathbf{x}
).
\]

A default rule may classify as \(1\) when

\[
\widehat{p}(\mathbf{x})>0.5.
\]

However, \(0.5\) is not universally optimal.

The threshold should reflect:

\[
\boxed{
\text{class prevalence},
}
\]

\[
\boxed{
\text{error costs},
}
\]

and

\[
\boxed{
\text{decision goals}.
}
\]

%%%%%%%%%%%%%%%%%%%%%%%%%%%%%%%%%%%%%%%%%%%%%%%%%%%%%%%%%%%%
\subsection{ROC Curve}
\label{subsec:roc-curve}
%%%%%%%%%%%%%%%%%%%%%%%%%%%%%%%%%%%%%%%%%%%%%%%%%%%%%%%%%%%%

A receiver operating characteristic, or ROC, curve plots

\[
\boxed{
\text{true-positive rate}
}
\]

against

\[
\boxed{
\text{false-positive rate}
=
1-\text{specificity}
}
\]

as the classification threshold varies.

%%%%%%%%%%%%%%%%%%%%%%%%%%%%%%%%%%%%%%%%%%%%%%%%%%%%%%%%%%%%
\subsection{Area Under the ROC Curve}
\label{subsec:auc}
%%%%%%%%%%%%%%%%%%%%%%%%%%%%%%%%%%%%%%%%%%%%%%%%%%%%%%%%%%%%

The area under the ROC curve is abbreviated AUC.

A probabilistic interpretation is approximately

\[
\boxed{
P
(
\text{a randomly chosen positive receives a larger score than a
randomly chosen negative}
).
}
\]

AUC evaluates ranking discrimination, not probability calibration.

%%%%%%%%%%%%%%%%%%%%%%%%%%%%%%%%%%%%%%%%%%%%%%%%%%%%%%%%%%%%
\subsection{Calibration}
\label{subsec:classification-calibration}
%%%%%%%%%%%%%%%%%%%%%%%%%%%%%%%%%%%%%%%%%%%%%%%%%%%%%%%%%%%%

A probabilistic classifier is calibrated if predicted probabilities agree
with observed frequencies.

For example, among observations assigned probability approximately

\[
0.8,
\]

roughly

\[
80\%
\]

should belong to the positive class.

A model can have strong discrimination and poor calibration.

%%%%%%%%%%%%%%%%%%%%%%%%%%%%%%%%%%%%%%%%%%%%%%%%%%%%%%%%%%%%
\subsection{Logarithmic Loss}
\label{subsec:log-loss}
%%%%%%%%%%%%%%%%%%%%%%%%%%%%%%%%%%%%%%%%%%%%%%%%%%%%%%%%%%%%

For binary response \(Y\in\{0,1\}\) and predicted probability \(p\),

\[
\boxed{
L
=
-
[
Y\log p
+
(1-Y)\log(1-p)
].
}
\]

This is binary log loss.

It heavily penalizes confident incorrect predictions.

%%%%%%%%%%%%%%%%%%%%%%%%%%%%%%%%%%%%%%%%%%%%%%%%%%%%%%%%%%%%
\subsection{Brier Score}
\label{subsec:brier-score-learning}
%%%%%%%%%%%%%%%%%%%%%%%%%%%%%%%%%%%%%%%%%%%%%%%%%%%%%%%%%%%%

For binary classification, the Brier score is

\[
\boxed{
\frac1n
\sum_{i=1}^{n}
(
Y_i-\widehat{p}_i
)^2.
}
\]

It evaluates probabilistic prediction and is sensitive to both calibration
and discrimination.

%%%%%%%%%%%%%%%%%%%%%%%%%%%%%%%%%%%%%%%%%%%%%%%%%%%%%%%%%%%%
\subsection{Logistic Regression for Classification}
\label{subsec:logistic-regression-learning}
%%%%%%%%%%%%%%%%%%%%%%%%%%%%%%%%%%%%%%%%%%%%%%%%%%%%%%%%%%%%

Logistic regression models

\[
\boxed{
P
(
Y=1
\mid
\mathbf{X}
)
=
p(\mathbf{X}),
}
\]

with

\[
\boxed{
\log
\left(
\frac{
p(\mathbf{X})
}{
1-p(\mathbf{X})
}
\right)
=
\beta_0
+
\mathbf{X}^T
\boldsymbol{\beta}.
}
\]

It is both an inferential model and a predictive classifier.

%%%%%%%%%%%%%%%%%%%%%%%%%%%%%%%%%%%%%%%%%%%%%%%%%%%%%%%%%%%%
\subsection{Logistic Decision Boundary}
\label{subsec:logistic-decision-boundary}
%%%%%%%%%%%%%%%%%%%%%%%%%%%%%%%%%%%%%%%%%%%%%%%%%%%%%%%%%%%%

Using threshold \(0.5\), classify as \(1\) when

\[
p(\mathbf{x})>0.5.
\]

Since

\[
\operatorname{logit}(0.5)=0,
\]

the boundary is

\[
\boxed{
\beta_0
+
\mathbf{x}^T
\boldsymbol{\beta}
=
0.
}
\]

Thus ordinary logistic regression produces a linear decision boundary in
predictor space.

%%%%%%%%%%%%%%%%%%%%%%%%%%%%%%%%%%%%%%%%%%%%%%%%%%%%%%%%%%%%
\subsection{\(k\)-Nearest Neighbors}
\label{subsec:k-nearest-neighbors}
%%%%%%%%%%%%%%%%%%%%%%%%%%%%%%%%%%%%%%%%%%%%%%%%%%%%%%%%%%%%

The \(k\)-nearest-neighbor method, abbreviated KNN, predicts using nearby
training observations.

For regression,

\[
\boxed{
\widehat{f}(\mathbf{x}_0)
=
\frac1k
\sum_{i\in N_k(\mathbf{x}_0)}
y_i,
}
\]

where

\[
N_k(\mathbf{x}_0)
\]

is the set of \(k\) nearest training observations.

%%%%%%%%%%%%%%%%%%%%%%%%%%%%%%%%%%%%%%%%%%%%%%%%%%%%%%%%%%%%
\subsection{KNN Classification}
\label{subsec:knn-classification}
%%%%%%%%%%%%%%%%%%%%%%%%%%%%%%%%%%%%%%%%%%%%%%%%%%%%%%%%%%%%

For classification, KNN predicts the majority class among the nearest
neighbors.

For binary data, an estimated probability is

\[
\boxed{
\widehat{p}(\mathbf{x}_0)
=
\frac{
\#\{
i\in N_k(\mathbf{x}_0):
Y_i=1
\}
}{
k
}.
}
\]

%%%%%%%%%%%%%%%%%%%%%%%%%%%%%%%%%%%%%%%%%%%%%%%%%%%%%%%%%%%%
\subsection{Choosing \(k\) in KNN}
\label{subsec:choosing-k-knn}
%%%%%%%%%%%%%%%%%%%%%%%%%%%%%%%%%%%%%%%%%%%%%%%%%%%%%%%%%%%%

Small \(k\):

\[
\boxed{
\text{high flexibility},
}
\]

\[
\boxed{
\text{low bias},
}
\]

\[
\boxed{
\text{high variance}.
}
\]

Large \(k\):

\[
\boxed{
\text{greater smoothing},
}
\]

\[
\boxed{
\text{higher bias},
}
\]

\[
\boxed{
\text{lower variance}.
}
\]

Cross-validation is commonly used to choose \(k\).

%%%%%%%%%%%%%%%%%%%%%%%%%%%%%%%%%%%%%%%%%%%%%%%%%%%%%%%%%%%%
\subsection{Scaling for Distance-Based Methods}
\label{subsec:scaling-distance-methods}
%%%%%%%%%%%%%%%%%%%%%%%%%%%%%%%%%%%%%%%%%%%%%%%%%%%%%%%%%%%%

KNN and clustering rely on distances.

If one predictor has much larger numerical scale than others, it can
dominate the distance calculation.

Standardization is therefore often essential.

%%%%%%%%%%%%%%%%%%%%%%%%%%%%%%%%%%%%%%%%%%%%%%%%%%%%%%%%%%%%
\subsection{Curse of Dimensionality}
\label{subsec:curse-dimensionality-learning}
%%%%%%%%%%%%%%%%%%%%%%%%%%%%%%%%%%%%%%%%%%%%%%%%%%%%%%%%%%%%

In high-dimensional spaces, observations become sparse.

Local neighborhoods contain less information unless sample size grows
rapidly.

Consequences include:

\[
\boxed{
\text{unstable nearest neighbors},
}
\]

\[
\boxed{
\text{high variance},
}
\]

and

\[
\boxed{
\text{distance concentration}.
}
\]

This is one aspect of the curse of dimensionality.

%%%%%%%%%%%%%%%%%%%%%%%%%%%%%%%%%%%%%%%%%%%%%%%%%%%%%%%%%%%%
\subsection{Decision Trees}
\label{subsec:decision-trees}
%%%%%%%%%%%%%%%%%%%%%%%%%%%%%%%%%%%%%%%%%%%%%%%%%%%%%%%%%%%%

Decision trees partition predictor space into regions.

A tree makes recursive binary splits such as

\[
\boxed{
X_j<c.
}
\]

Each terminal region receives a prediction.

Trees can model nonlinearities and interactions automatically.

%%%%%%%%%%%%%%%%%%%%%%%%%%%%%%%%%%%%%%%%%%%%%%%%%%%%%%%%%%%%
\subsection{Regression Trees}
\label{subsec:regression-trees}
%%%%%%%%%%%%%%%%%%%%%%%%%%%%%%%%%%%%%%%%%%%%%%%%%%%%%%%%%%%%

For a regression tree, each terminal node predicts the mean response of
training observations in that region.

Suppose terminal region \(R_m\) contains observations \(i\in R_m\).

Then

\[
\boxed{
\widehat{c}_m
=
\frac1{|R_m|}
\sum_{i\in R_m}
y_i.
}
\]

%%%%%%%%%%%%%%%%%%%%%%%%%%%%%%%%%%%%%%%%%%%%%%%%%%%%%%%%%%%%
\subsection{Regression Tree Split Criterion}
\label{subsec:regression-tree-split}
%%%%%%%%%%%%%%%%%%%%%%%%%%%%%%%%%%%%%%%%%%%%%%%%%%%%%%%%%%%%

A split is often chosen to minimize residual sum of squares:

\[
\boxed{
\sum_{i\in R_1}
(
y_i-\overline{y}_{R_1}
)^2
+
\sum_{i\in R_2}
(
y_i-\overline{y}_{R_2}
)^2.
}
\]

The best predictor and split point are selected greedily.

%%%%%%%%%%%%%%%%%%%%%%%%%%%%%%%%%%%%%%%%%%%%%%%%%%%%%%%%%%%%
\subsection{Classification Trees}
\label{subsec:classification-trees}
%%%%%%%%%%%%%%%%%%%%%%%%%%%%%%%%%%%%%%%%%%%%%%%%%%%%%%%%%%%%

A classification tree assigns a class to each terminal region.

If

\[
\widehat{p}_{mk}
\]

is the proportion of observations in region \(m\) belonging to class
\(k\), the predicted class is commonly

\[
\boxed{
\arg\max_k
\widehat{p}_{mk}.
}
\]

%%%%%%%%%%%%%%%%%%%%%%%%%%%%%%%%%%%%%%%%%%%%%%%%%%%%%%%%%%%%
\subsection{Gini Impurity}
\label{subsec:gini-impurity}
%%%%%%%%%%%%%%%%%%%%%%%%%%%%%%%%%%%%%%%%%%%%%%%%%%%%%%%%%%%%

For node \(m\), Gini impurity is

\[
\boxed{
G_m
=
\sum_{k=1}^{K}
\widehat{p}_{mk}
(
1-\widehat{p}_{mk}
).
}
\]

Equivalently,

\[
\boxed{
G_m
=
1
-
\sum_{k=1}^{K}
\widehat{p}_{mk}^2.
}
\]

It is small when one class dominates the node.

%%%%%%%%%%%%%%%%%%%%%%%%%%%%%%%%%%%%%%%%%%%%%%%%%%%%%%%%%%%%
\subsection{Classification Entropy}
\label{subsec:classification-entropy}
%%%%%%%%%%%%%%%%%%%%%%%%%%%%%%%%%%%%%%%%%%%%%%%%%%%%%%%%%%%%

Another impurity measure is entropy:

\[
\boxed{
H_m
=
-
\sum_{k=1}^{K}
\widehat{p}_{mk}
\log
\widehat{p}_{mk}.
}
\]

Entropy is zero when the node contains observations from only one class.

%%%%%%%%%%%%%%%%%%%%%%%%%%%%%%%%%%%%%%%%%%%%%%%%%%%%%%%%%%%%
\subsection{Tree Overfitting}
\label{subsec:tree-overfitting}
%%%%%%%%%%%%%%%%%%%%%%%%%%%%%%%%%%%%%%%%%%%%%%%%%%%%%%%%%%%%

A fully grown tree can create very small terminal nodes and fit training
noise.

Thus trees can have high variance.

Restricting tree complexity or pruning can improve generalization.

%%%%%%%%%%%%%%%%%%%%%%%%%%%%%%%%%%%%%%%%%%%%%%%%%%%%%%%%%%%%
\subsection{Cost-Complexity Pruning}
\label{subsec:cost-complexity-pruning}
%%%%%%%%%%%%%%%%%%%%%%%%%%%%%%%%%%%%%%%%%%%%%%%%%%%%%%%%%%%%

A pruned regression tree may minimize an objective such as

\[
\boxed{
RSS(T)
+
\alpha
|T|,
}
\]

where

\[
|T|
\]

is the number of terminal nodes.

The parameter \(\alpha\) controls the penalty on tree size.

Cross-validation can choose \(\alpha\).

%%%%%%%%%%%%%%%%%%%%%%%%%%%%%%%%%%%%%%%%%%%%%%%%%%%%%%%%%%%%
\subsection{Advantages of Trees}
\label{subsec:advantages-trees}
%%%%%%%%%%%%%%%%%%%%%%%%%%%%%%%%%%%%%%%%%%%%%%%%%%%%%%%%%%%%

Trees are attractive because they:

\[
\boxed{
\text{handle nonlinearities},
}
\]

\[
\boxed{
\text{capture interactions},
}
\]

\[
\boxed{
\text{require little functional-form specification},
}
\]

and can be

\[
\boxed{
\text{visually interpretable}.
}
\]

%%%%%%%%%%%%%%%%%%%%%%%%%%%%%%%%%%%%%%%%%%%%%%%%%%%%%%%%%%%%
\subsection{Limitations of Single Trees}
\label{subsec:limitations-single-trees}
%%%%%%%%%%%%%%%%%%%%%%%%%%%%%%%%%%%%%%%%%%%%%%%%%%%%%%%%%%%%

Single trees can be:

\[
\boxed{
\text{unstable},
}
\]

\[
\boxed{
\text{high variance},
}
\]

and less accurate than ensemble methods.

Small changes in training data can produce very different tree
structures.

%%%%%%%%%%%%%%%%%%%%%%%%%%%%%%%%%%%%%%%%%%%%%%%%%%%%%%%%%%%%
\subsection{Bootstrap Aggregation}
\label{subsec:bagging}
%%%%%%%%%%%%%%%%%%%%%%%%%%%%%%%%%%%%%%%%%%%%%%%%%%%%%%%%%%%%

Bootstrap aggregation, or bagging, reduces variance by averaging models
fit to bootstrap samples.

For \(B\) bootstrap models,

\[
\widehat{f}^{(1)},
\ldots,
\widehat{f}^{(B)},
\]

the bagged regression predictor is

\[
\boxed{
\widehat{f}_{bag}(\mathbf{x})
=
\frac1B
\sum_{b=1}^{B}
\widehat{f}^{(b)}(\mathbf{x}).
}
\]

%%%%%%%%%%%%%%%%%%%%%%%%%%%%%%%%%%%%%%%%%%%%%%%%%%%%%%%%%%%%
\subsection{Why Bagging Reduces Variance}
\label{subsec:why-bagging-reduces-variance}
%%%%%%%%%%%%%%%%%%%%%%%%%%%%%%%%%%%%%%%%%%%%%%%%%%%%%%%%%%%%

Suppose \(B\) predictors have common variance \(\sigma^2\) and pairwise
correlation \(\rho\).

The variance of their average is approximately

\[
\boxed{
\rho\sigma^2
+
\frac{
1-\rho
}{
B
}
\sigma^2.
}
\]

Averaging is most effective when individual predictions are not too
strongly correlated.

%%%%%%%%%%%%%%%%%%%%%%%%%%%%%%%%%%%%%%%%%%%%%%%%%%%%%%%%%%%%
\subsection{Random Forests}
\label{subsec:random-forests}
%%%%%%%%%%%%%%%%%%%%%%%%%%%%%%%%%%%%%%%%%%%%%%%%%%%%%%%%%%%%

Random forests extend bagging by randomly restricting candidate
predictors at each tree split.

This decorrelates the trees.

The final prediction averages many trees.

For regression,

\[
\boxed{
\widehat{f}_{RF}(\mathbf{x})
=
\frac1B
\sum_{b=1}^{B}
\widehat{f}_b(\mathbf{x}).
}
\]

%%%%%%%%%%%%%%%%%%%%%%%%%%%%%%%%%%%%%%%%%%%%%%%%%%%%%%%%%%%%
\subsection{Random-Forest Classification}
\label{subsec:random-forest-classification}
%%%%%%%%%%%%%%%%%%%%%%%%%%%%%%%%%%%%%%%%%%%%%%%%%%%%%%%%%%%%

For classification, each tree votes for a class.

The forest may predict using:

\[
\boxed{
\text{majority vote},
}
\]

or average estimated class probabilities across trees.

%%%%%%%%%%%%%%%%%%%%%%%%%%%%%%%%%%%%%%%%%%%%%%%%%%%%%%%%%%%%
\subsection{Out-of-Bag Observations}
\label{subsec:out-of-bag}
%%%%%%%%%%%%%%%%%%%%%%%%%%%%%%%%%%%%%%%%%%%%%%%%%%%%%%%%%%%%

A bootstrap sample leaves out some training observations.

These are called out-of-bag, or OOB, observations.

For each observation, predictions can be made using only trees whose
bootstrap samples did not contain that observation.

This provides an internal estimate of prediction error.

%%%%%%%%%%%%%%%%%%%%%%%%%%%%%%%%%%%%%%%%%%%%%%%%%%%%%%%%%%%%
\subsection{Variable Importance}
\label{subsec:variable-importance}
%%%%%%%%%%%%%%%%%%%%%%%%%%%%%%%%%%%%%%%%%%%%%%%%%%%%%%%%%%%%

Ensemble methods often report variable-importance measures.

Examples include:

\[
\boxed{
\text{split-improvement importance},
}
\]

and

\[
\boxed{
\text{permutation importance}.
}
\]

Importance measures quantify predictive contribution under a fitted model
but do not automatically imply causality.

%%%%%%%%%%%%%%%%%%%%%%%%%%%%%%%%%%%%%%%%%%%%%%%%%%%%%%%%%%%%
\subsection{Permutation Importance}
\label{subsec:permutation-importance}
%%%%%%%%%%%%%%%%%%%%%%%%%%%%%%%%%%%%%%%%%%%%%%%%%%%%%%%%%%%%

Permutation importance randomly permutes one predictor in validation data
and measures the resulting increase in prediction error.

If permuting \(X_j\) greatly harms performance, the model relies strongly
on \(X_j\).

Correlated predictors can complicate interpretation.

%%%%%%%%%%%%%%%%%%%%%%%%%%%%%%%%%%%%%%%%%%%%%%%%%%%%%%%%%%%%
\subsection{Boosting}
\label{subsec:boosting}
%%%%%%%%%%%%%%%%%%%%%%%%%%%%%%%%%%%%%%%%%%%%%%%%%%%%%%%%%%%%

Boosting builds models sequentially.

Each new learner attempts to improve upon the current ensemble.

A generic additive model is

\[
\boxed{
F_M(\mathbf{x})
=
\sum_{m=1}^{M}
\nu_m
g_m(\mathbf{x}),
}
\]

where the \(g_m\) are weak learners.

%%%%%%%%%%%%%%%%%%%%%%%%%%%%%%%%%%%%%%%%%%%%%%%%%%%%%%%%%%%%
\subsection{Gradient Boosting}
\label{subsec:gradient-boosting}
%%%%%%%%%%%%%%%%%%%%%%%%%%%%%%%%%%%%%%%%%%%%%%%%%%%%%%%%%%%%

Gradient boosting interprets boosting as numerical optimization in
function space.

At each step, a new learner is fit approximately to the negative gradient
of the loss with respect to current predictions.

For squared-error loss, this corresponds to fitting residuals.

%%%%%%%%%%%%%%%%%%%%%%%%%%%%%%%%%%%%%%%%%%%%%%%%%%%%%%%%%%%%
\subsection{Learning Rate}
\label{subsec:boosting-learning-rate}
%%%%%%%%%%%%%%%%%%%%%%%%%%%%%%%%%%%%%%%%%%%%%%%%%%%%%%%%%%%%

A learning rate

\[
\boxed{
0<\nu\leq1
}
\]

controls how much each new learner contributes.

Small learning rates generally require more boosting iterations but can
improve generalization.

The number of iterations and learning rate should be tuned jointly.

%%%%%%%%%%%%%%%%%%%%%%%%%%%%%%%%%%%%%%%%%%%%%%%%%%%%%%%%%%%%
\subsection{Bagging Versus Boosting}
\label{subsec:bagging-versus-boosting}
%%%%%%%%%%%%%%%%%%%%%%%%%%%%%%%%%%%%%%%%%%%%%%%%%%%%%%%%%%%%

Bagging builds many models largely in parallel and averages them.

Boosting builds models sequentially, with each learner responding to the
current ensemble's errors.

Thus:

\[
\boxed{
\text{bagging}
\rightarrow
\text{variance reduction},
}
\]

while

\[
\boxed{
\text{boosting}
\rightarrow
\text{iterative bias and error reduction}.
}
\]

This distinction is conceptual rather than absolute.

%%%%%%%%%%%%%%%%%%%%%%%%%%%%%%%%%%%%%%%%%%%%%%%%%%%%%%%%%%%%
\subsection{Support-Vector Classifiers}
\label{subsec:support-vector-classifiers}
%%%%%%%%%%%%%%%%%%%%%%%%%%%%%%%%%%%%%%%%%%%%%%%%%%%%%%%%%%%%

Suppose two classes are linearly separable.

A separating hyperplane has form

\[
\boxed{
\beta_0
+
\boldsymbol{\beta}^T
\mathbf{x}
=
0.
}
\]

A support-vector classifier seeks a separating boundary with a large
margin between the classes.

%%%%%%%%%%%%%%%%%%%%%%%%%%%%%%%%%%%%%%%%%%%%%%%%%%%%%%%%%%%%
\subsection{Margin}
\label{subsec:svm-margin}
%%%%%%%%%%%%%%%%%%%%%%%%%%%%%%%%%%%%%%%%%%%%%%%%%%%%%%%%%%%%

The margin is the distance between the decision boundary and the nearest
training observations.

The observations determining the margin are called support vectors.

A larger margin can improve robustness to small perturbations.

%%%%%%%%%%%%%%%%%%%%%%%%%%%%%%%%%%%%%%%%%%%%%%%%%%%%%%%%%%%%
\subsection{Soft-Margin Classification}
\label{subsec:soft-margin}
%%%%%%%%%%%%%%%%%%%%%%%%%%%%%%%%%%%%%%%%%%%%%%%%%%%%%%%%%%%%

Perfect separation may be impossible or undesirable.

Soft-margin methods allow some violations through slack variables

\[
\xi_i\geq0.
\]

A typical optimization has the form

\[
\boxed{
\min_{\boldsymbol{\beta},\xi}
\frac12
\|
\boldsymbol{\beta}
\|_2^2
+
C
\sum_i
\xi_i.
}
\]

The parameter \(C\) controls the penalty for violations.

%%%%%%%%%%%%%%%%%%%%%%%%%%%%%%%%%%%%%%%%%%%%%%%%%%%%%%%%%%%%
\subsection{Kernel Trick}
\label{subsec:kernel-trick}
%%%%%%%%%%%%%%%%%%%%%%%%%%%%%%%%%%%%%%%%%%%%%%%%%%%%%%%%%%%%

Some algorithms depend only on inner products

\[
\mathbf{x}_i^T
\mathbf{x}_j.
\]

A kernel replaces these with

\[
\boxed{
K
(
\mathbf{x}_i,
\mathbf{x}_j
)
=
\langle
\phi(\mathbf{x}_i),
\phi(\mathbf{x}_j)
\rangle,
}
\]

where \(\phi\) maps inputs to a possibly high-dimensional feature space.

The mapping need not be computed explicitly.

%%%%%%%%%%%%%%%%%%%%%%%%%%%%%%%%%%%%%%%%%%%%%%%%%%%%%%%%%%%%
\subsection{Common Kernels}
\label{subsec:common-kernels}
%%%%%%%%%%%%%%%%%%%%%%%%%%%%%%%%%%%%%%%%%%%%%%%%%%%%%%%%%%%%

A linear kernel is

\[
\boxed{
K(\mathbf{x},\mathbf{z})
=
\mathbf{x}^T\mathbf{z}.
}
\]

A polynomial kernel is

\[
\boxed{
K(\mathbf{x},\mathbf{z})
=
(
c+\mathbf{x}^T\mathbf{z}
)^d.
}
\]

A radial basis function kernel is

\[
\boxed{
K(\mathbf{x},\mathbf{z})
=
\exp
\left[
-\gamma
\|
\mathbf{x}-\mathbf{z}
\|_2^2
\right].
}
\]

%%%%%%%%%%%%%%%%%%%%%%%%%%%%%%%%%%%%%%%%%%%%%%%%%%%%%%%%%%%%
\subsection{Kernel Validity}
\label{subsec:kernel-validity}
%%%%%%%%%%%%%%%%%%%%%%%%%%%%%%%%%%%%%%%%%%%%%%%%%%%%%%%%%%%%

A valid kernel must generate a positive semidefinite Gram matrix.

For training observations,

\[
\boxed{
K_{ij}
=
K
(
\mathbf{x}_i,
\mathbf{x}_j
).
}
\]

Then

\[
\boxed{
\mathbf{a}^TK\mathbf{a}
\geq0
}
\]

for every vector \(\mathbf{a}\).

This connects kernel methods to positive semidefinite matrix theory.

%%%%%%%%%%%%%%%%%%%%%%%%%%%%%%%%%%%%%%%%%%%%%%%%%%%%%%%%%%%%
\subsection{Support-Vector Machines}
\label{subsec:support-vector-machines}
%%%%%%%%%%%%%%%%%%%%%%%%%%%%%%%%%%%%%%%%%%%%%%%%%%%%%%%%%%%%

A support-vector machine combines margin-based classification with a
kernel.

The resulting decision function has a form such as

\[
\boxed{
f(\mathbf{x})
=
\beta_0
+
\sum_{i=1}^{n}
\alpha_i
K
(
\mathbf{x}_i,
\mathbf{x}
).
}
\]

Only training observations with nonzero \(\alpha_i\) act as support
vectors.

%%%%%%%%%%%%%%%%%%%%%%%%%%%%%%%%%%%%%%%%%%%%%%%%%%%%%%%%%%%%
\subsection{Principal Components Regression}
\label{subsec:principal-components-regression}
%%%%%%%%%%%%%%%%%%%%%%%%%%%%%%%%%%%%%%%%%%%%%%%%%%%%%%%%%%%%

Principal components regression, abbreviated PCR, first performs PCA on
the predictors.

Then the response is regressed on the first \(m\) principal-component
scores.

Thus

\[
\boxed{
X
\longrightarrow
\text{PCA scores}
\longrightarrow
\text{regression}.
}
\]

%%%%%%%%%%%%%%%%%%%%%%%%%%%%%%%%%%%%%%%%%%%%%%%%%%%%%%%%%%%%
\subsection{Limitation of PCR}
\label{subsec:pcr-limitation}
%%%%%%%%%%%%%%%%%%%%%%%%%%%%%%%%%%%%%%%%%%%%%%%%%%%%%%%%%%%%

PCA chooses directions with large predictor variance.

It does not use \(Y\).

Therefore a low-variance predictor direction may be highly predictive yet
discarded by PCR.

This is why unsupervised dimension reduction is not always optimal for
prediction.

%%%%%%%%%%%%%%%%%%%%%%%%%%%%%%%%%%%%%%%%%%%%%%%%%%%%%%%%%%%%
\subsection{Partial Least Squares Preview}
\label{subsec:partial-least-squares}
%%%%%%%%%%%%%%%%%%%%%%%%%%%%%%%%%%%%%%%%%%%%%%%%%%%%%%%%%%%%

Partial least squares, abbreviated PLS, constructs components using both:

\[
\boxed{
\text{predictor variation}
}
\]

and

\[
\boxed{
\text{association with the response}.
}
\]

Thus PLS is a supervised dimension-reduction method.

%%%%%%%%%%%%%%%%%%%%%%%%%%%%%%%%%%%%%%%%%%%%%%%%%%%%%%%%%%%%
\subsection{High-Dimensional Learning}
\label{subsec:high-dimensional-learning}
%%%%%%%%%%%%%%%%%%%%%%%%%%%%%%%%%%%%%%%%%%%%%%%%%%%%%%%%%%%%

High-dimensional settings occur when

\[
\boxed{
p
\text{ is large relative to }
n,
}
\]

or even

\[
\boxed{
p>n.
}
\]

Ordinary least squares may then be unstable or nonunique.

Regularization and dimension reduction become essential.

%%%%%%%%%%%%%%%%%%%%%%%%%%%%%%%%%%%%%%%%%%%%%%%%%%%%%%%%%%%%
\subsection{Sparsity}
\label{subsec:sparsity}
%%%%%%%%%%%%%%%%%%%%%%%%%%%%%%%%%%%%%%%%%%%%%%%%%%%%%%%%%%%%

A sparse model assumes only a relatively small number of predictors have
substantial effects.

If

\[
\boldsymbol{\beta}
=
(
\beta_1,\ldots,\beta_p
),
\]

sparsity means many

\[
\boxed{
\beta_j=0.
}
\]

Lasso and related methods explicitly encourage sparsity.

%%%%%%%%%%%%%%%%%%%%%%%%%%%%%%%%%%%%%%%%%%%%%%%%%%%%%%%%%%%%
\subsection{Unsupervised Clustering}
\label{subsec:unsupervised-clustering}
%%%%%%%%%%%%%%%%%%%%%%%%%%%%%%%%%%%%%%%%%%%%%%%%%%%%%%%%%%%%

Clustering attempts to divide observations into groups with:

\[
\boxed{
\text{high within-group similarity}
}
\]

and

\[
\boxed{
\text{high between-group dissimilarity}.
}
\]

Unlike classification, true class labels are not supplied.

%%%%%%%%%%%%%%%%%%%%%%%%%%%%%%%%%%%%%%%%%%%%%%%%%%%%%%%%%%%%
\subsection{\(k\)-Means Clustering}
\label{subsec:kmeans-clustering}
%%%%%%%%%%%%%%%%%%%%%%%%%%%%%%%%%%%%%%%%%%%%%%%%%%%%%%%%%%%%

Given \(K\) clusters

\[
C_1,\ldots,C_K,
\]

\(k\)-means minimizes

\[
\boxed{
\sum_{k=1}^{K}
\sum_{i\in C_k}
\|
\mathbf{x}_i
-
\boldsymbol{\mu}_k
\|_2^2,
}
\]

where

\[
\boldsymbol{\mu}_k
\]

is the mean vector of cluster \(k\).

%%%%%%%%%%%%%%%%%%%%%%%%%%%%%%%%%%%%%%%%%%%%%%%%%%%%%%%%%%%%
\subsection{\(k\)-Means Algorithm}
\label{subsec:kmeans-algorithm}
%%%%%%%%%%%%%%%%%%%%%%%%%%%%%%%%%%%%%%%%%%%%%%%%%%%%%%%%%%%%

A common iterative algorithm is:

\begin{enumerate}
    \item choose \(K\) initial cluster centers;
    \item assign each observation to its nearest center;
    \item recompute each cluster center;
    \item reassign observations;
    \item repeat until assignments stabilize or the objective stops
          improving.
\end{enumerate}

The algorithm decreases the within-cluster objective at each iteration.

%%%%%%%%%%%%%%%%%%%%%%%%%%%%%%%%%%%%%%%%%%%%%%%%%%%%%%%%%%%%
\subsection{\(k\)-Means Local Optima}
\label{subsec:kmeans-local-optima}
%%%%%%%%%%%%%%%%%%%%%%%%%%%%%%%%%%%%%%%%%%%%%%%%%%%%%%%%%%%%

The \(k\)-means objective is not generally convex in cluster assignments
and centers jointly.

Different initializations may produce different solutions.

Therefore multiple random starts are commonly used.

%%%%%%%%%%%%%%%%%%%%%%%%%%%%%%%%%%%%%%%%%%%%%%%%%%%%%%%%%%%%
\subsection{Choosing the Number of Clusters}
\label{subsec:choosing-number-clusters}
%%%%%%%%%%%%%%%%%%%%%%%%%%%%%%%%%%%%%%%%%%%%%%%%%%%%%%%%%%%%

Possible approaches include:

\[
\boxed{
\text{within-cluster variation},
}
\]

\[
\boxed{
\text{silhouette measures},
}
\]

\[
\boxed{
\text{stability},
}
\]

and

\[
\boxed{
\text{scientific interpretability}.
}
\]

No universal rule identifies the true number of clusters.

%%%%%%%%%%%%%%%%%%%%%%%%%%%%%%%%%%%%%%%%%%%%%%%%%%%%%%%%%%%%
\subsection{Silhouette Concept}
\label{subsec:silhouette}
%%%%%%%%%%%%%%%%%%%%%%%%%%%%%%%%%%%%%%%%%%%%%%%%%%%%%%%%%%%%

For observation \(i\), let

\[
a(i)
\]

be its average distance to points in its own cluster.

Let

\[
b(i)
\]

be the smallest average distance to another cluster.

The silhouette value is

\[
\boxed{
s(i)
=
\frac{
b(i)-a(i)
}{
\max\{a(i),b(i)\}
}.
}
\]

Values near \(1\) indicate stronger separation.

%%%%%%%%%%%%%%%%%%%%%%%%%%%%%%%%%%%%%%%%%%%%%%%%%%%%%%%%%%%%
\subsection{Hierarchical Clustering}
\label{subsec:hierarchical-clustering}
%%%%%%%%%%%%%%%%%%%%%%%%%%%%%%%%%%%%%%%%%%%%%%%%%%%%%%%%%%%%

Hierarchical clustering creates a nested sequence of clusters.

Agglomerative clustering begins with each observation as its own cluster
and repeatedly merges groups.

The result is represented by a dendrogram.

%%%%%%%%%%%%%%%%%%%%%%%%%%%%%%%%%%%%%%%%%%%%%%%%%%%%%%%%%%%%
\subsection{Linkage Criteria}
\label{subsec:linkage-criteria}
%%%%%%%%%%%%%%%%%%%%%%%%%%%%%%%%%%%%%%%%%%%%%%%%%%%%%%%%%%%%

To merge clusters, a distance between clusters must be defined.

Common choices include:

\[
\boxed{
\text{single linkage},
}
\]

\[
\boxed{
\text{complete linkage},
}
\]

\[
\boxed{
\text{average linkage},
}
\]

and

\[
\boxed{
\text{Ward-type criteria}.
}
\]

Different linkage choices can produce substantially different clusters.

%%%%%%%%%%%%%%%%%%%%%%%%%%%%%%%%%%%%%%%%%%%%%%%%%%%%%%%%%%%%
\subsection{Single Linkage}
\label{subsec:single-linkage}
%%%%%%%%%%%%%%%%%%%%%%%%%%%%%%%%%%%%%%%%%%%%%%%%%%%%%%%%%%%%

Single linkage defines cluster distance using the closest pair:

\[
\boxed{
d(A,B)
=
\min_{
i\in A,
j\in B
}
d(i,j).
}
\]

It can identify elongated clusters but may produce chaining.

%%%%%%%%%%%%%%%%%%%%%%%%%%%%%%%%%%%%%%%%%%%%%%%%%%%%%%%%%%%%
\subsection{Complete Linkage}
\label{subsec:complete-linkage}
%%%%%%%%%%%%%%%%%%%%%%%%%%%%%%%%%%%%%%%%%%%%%%%%%%%%%%%%%%%%

Complete linkage uses the farthest pair:

\[
\boxed{
d(A,B)
=
\max_{
i\in A,
j\in B
}
d(i,j).
}
\]

It tends to favor more compact clusters.

%%%%%%%%%%%%%%%%%%%%%%%%%%%%%%%%%%%%%%%%%%%%%%%%%%%%%%%%%%%%
\subsection{Average Linkage}
\label{subsec:average-linkage}
%%%%%%%%%%%%%%%%%%%%%%%%%%%%%%%%%%%%%%%%%%%%%%%%%%%%%%%%%%%%

Average linkage uses the average pairwise distance:

\[
\boxed{
d(A,B)
=
\frac1{
|A||B|
}
\sum_{i\in A}
\sum_{j\in B}
d(i,j).
}
\]

It balances characteristics of single and complete linkage.

%%%%%%%%%%%%%%%%%%%%%%%%%%%%%%%%%%%%%%%%%%%%%%%%%%%%%%%%%%%%
\subsection{Distance Measures}
\label{subsec:distance-measures-learning}
%%%%%%%%%%%%%%%%%%%%%%%%%%%%%%%%%%%%%%%%%%%%%%%%%%%%%%%%%%%%

Euclidean distance is

\[
\boxed{
d(\mathbf{x},\mathbf{z})
=
\sqrt{
\sum_{j=1}^{p}
(
x_j-z_j
)^2
}.
}
\]

Manhattan distance is

\[
\boxed{
d_1
(
\mathbf{x},
\mathbf{z}
)
=
\sum_{j=1}^{p}
|
x_j-z_j
|.
}
\]

The chosen distance defines what ``similar'' means.

%%%%%%%%%%%%%%%%%%%%%%%%%%%%%%%%%%%%%%%%%%%%%%%%%%%%%%%%%%%%
\subsection{Mahalanobis Distance in Learning}
\label{subsec:mahalanobis-learning}
%%%%%%%%%%%%%%%%%%%%%%%%%%%%%%%%%%%%%%%%%%%%%%%%%%%%%%%%%%%%

Mahalanobis distance is

\[
\boxed{
d_M
(
\mathbf{x},
\mathbf{z}
)^2
=
(
\mathbf{x}-\mathbf{z}
)^T
\Sigma^{-1}
(
\mathbf{x}-\mathbf{z}
).
}
\]

It accounts for:

\[
\boxed{
\text{predictor scale}
}
\]

and

\[
\boxed{
\text{correlation}.
}
\]

A stable covariance estimate is required.

%%%%%%%%%%%%%%%%%%%%%%%%%%%%%%%%%%%%%%%%%%%%%%%%%%%%%%%%%%%%
\subsection{Ensemble Learning}
\label{subsec:ensemble-learning}
%%%%%%%%%%%%%%%%%%%%%%%%%%%%%%%%%%%%%%%%%%%%%%%%%%%%%%%%%%%%

An ensemble combines several models.

The objective is to obtain better prediction than a typical individual
model.

Examples include:

\[
\boxed{
\text{bagging},
}
\]

\[
\boxed{
\text{random forests},
}
\]

\[
\boxed{
\text{boosting},
}
\]

and

\[
\boxed{
\text{stacking}.
}
\]

%%%%%%%%%%%%%%%%%%%%%%%%%%%%%%%%%%%%%%%%%%%%%%%%%%%%%%%%%%%%
\subsection{Stacking Preview}
\label{subsec:stacking-preview}
%%%%%%%%%%%%%%%%%%%%%%%%%%%%%%%%%%%%%%%%%%%%%%%%%%%%%%%%%%%%

Stacking combines predictions from several base models.

Suppose

\[
\widehat{f}_1,\ldots,\widehat{f}_M
\]

are base predictors.

A meta-model learns how to combine

\[
\boxed{
\widehat{f}_1(\mathbf{x}),
\ldots,
\widehat{f}_M(\mathbf{x})
}
\]

using out-of-fold predictions.

This can adaptively combine models with different strengths.

%%%%%%%%%%%%%%%%%%%%%%%%%%%%%%%%%%%%%%%%%%%%%%%%%%%%%%%%%%%%
\subsection{Hyperparameters}
\label{subsec:hyperparameters-learning}
%%%%%%%%%%%%%%%%%%%%%%%%%%%%%%%%%%%%%%%%%%%%%%%%%%%%%%%%%%%%

A hyperparameter controls model complexity or fitting behavior but is not
typically estimated directly by ordinary likelihood fitting.

Examples include:

\[
\boxed{
\lambda
\text{ in ridge regression},
}
\]

\[
\boxed{
k
\text{ in KNN},
}
\]

\[
\boxed{
\text{tree depth},
}
\]

and

\[
\boxed{
\text{learning rate}.
}
\]

Hyperparameters are commonly selected through validation.

%%%%%%%%%%%%%%%%%%%%%%%%%%%%%%%%%%%%%%%%%%%%%%%%%%%%%%%%%%%%
\subsection{Grid Search}
\label{subsec:grid-search}
%%%%%%%%%%%%%%%%%%%%%%%%%%%%%%%%%%%%%%%%%%%%%%%%%%%%%%%%%%%%

Grid search evaluates a predetermined set of hyperparameter combinations.

For example,

\[
\lambda
\in
\{
0.001,
0.01,
0.1,
1,
10
\}.
\]

The combination with best validation performance is selected.

Grid search becomes expensive when many hyperparameters are tuned.

%%%%%%%%%%%%%%%%%%%%%%%%%%%%%%%%%%%%%%%%%%%%%%%%%%%%%%%%%%%%
\subsection{Random Search Preview}
\label{subsec:random-search}
%%%%%%%%%%%%%%%%%%%%%%%%%%%%%%%%%%%%%%%%%%%%%%%%%%%%%%%%%%%%

Random search samples hyperparameter combinations from specified
distributions.

It can explore high-dimensional tuning spaces more efficiently than a
dense Cartesian grid when only a few hyperparameters strongly affect
performance.

%%%%%%%%%%%%%%%%%%%%%%%%%%%%%%%%%%%%%%%%%%%%%%%%%%%%%%%%%%%%
\subsection{Model Selection Versus Model Assessment}
\label{subsec:model-selection-versus-assessment}
%%%%%%%%%%%%%%%%%%%%%%%%%%%%%%%%%%%%%%%%%%%%%%%%%%%%%%%%%%%%

Model selection chooses among candidate procedures.

Model assessment estimates the performance of the selected procedure.

Using the same data repeatedly for both purposes creates optimistic bias.

This is why independent test sets or nested resampling are important.

%%%%%%%%%%%%%%%%%%%%%%%%%%%%%%%%%%%%%%%%%%%%%%%%%%%%%%%%%%%%
\subsection{Learning Curves}
\label{subsec:learning-curves}
%%%%%%%%%%%%%%%%%%%%%%%%%%%%%%%%%%%%%%%%%%%%%%%%%%%%%%%%%%%%

A learning curve plots prediction performance against training-set size.

It can help diagnose whether additional data may improve performance.

For example:

\[
\boxed{
\text{large train-test gap}
\rightarrow
\text{possible high variance},
}
\]

while

\[
\boxed{
\text{both errors high}
\rightarrow
\text{possible high bias}.
}
\]

%%%%%%%%%%%%%%%%%%%%%%%%%%%%%%%%%%%%%%%%%%%%%%%%%%%%%%%%%%%%
\subsection{Imbalanced Classification}
\label{subsec:imbalanced-classification}
%%%%%%%%%%%%%%%%%%%%%%%%%%%%%%%%%%%%%%%%%%%%%%%%%%%%%%%%%%%%

Suppose only

\[
1\%
\]

of observations belong to the positive class.

A classifier predicting the negative class for everyone achieves

\[
99\%
\]

accuracy but has no ability to detect positives.

Thus accuracy alone can be misleading.

%%%%%%%%%%%%%%%%%%%%%%%%%%%%%%%%%%%%%%%%%%%%%%%%%%%%%%%%%%%%
\subsection{Metrics for Imbalanced Data}
\label{subsec:imbalanced-metrics}
%%%%%%%%%%%%%%%%%%%%%%%%%%%%%%%%%%%%%%%%%%%%%%%%%%%%%%%%%%%%

Useful metrics include:

\[
\boxed{
\text{precision},
}
\]

\[
\boxed{
\text{recall},
}
\]

\[
\boxed{
F_1,
}
\]

\[
\boxed{
\text{precision--recall curves},
}
\]

and cost-sensitive measures.

The best metric depends on the decision problem.

%%%%%%%%%%%%%%%%%%%%%%%%%%%%%%%%%%%%%%%%%%%%%%%%%%%%%%%%%%%%
\subsection{Class Weighting}
\label{subsec:class-weighting}
%%%%%%%%%%%%%%%%%%%%%%%%%%%%%%%%%%%%%%%%%%%%%%%%%%%%%%%%%%%%

A classification loss can give different weights to observations from
different classes.

For example,

\[
\boxed{
L
=
w_1L_1
+
w_0L_0.
}
\]

Larger positive-class weight can make false negatives more costly.

Weights should reflect the intended objective rather than merely forcing
balanced predictions.

%%%%%%%%%%%%%%%%%%%%%%%%%%%%%%%%%%%%%%%%%%%%%%%%%%%%%%%%%%%%
\subsection{Resampling for Class Imbalance}
\label{subsec:resampling-class-imbalance}
%%%%%%%%%%%%%%%%%%%%%%%%%%%%%%%%%%%%%%%%%%%%%%%%%%%%%%%%%%%%

Training data may be rebalanced through:

\[
\boxed{
\text{oversampling},
}
\]

or

\[
\boxed{
\text{undersampling}.
}
\]

Any such resampling must occur inside training folds to avoid leakage.

The final predicted probabilities may require calibration to the true
class prevalence.

%%%%%%%%%%%%%%%%%%%%%%%%%%%%%%%%%%%%%%%%%%%%%%%%%%%%%%%%%%%%
\subsection{Interpretability}
\label{subsec:model-interpretability}
%%%%%%%%%%%%%%%%%%%%%%%%%%%%%%%%%%%%%%%%%%%%%%%%%%%%%%%%%%%%

Interpretability concerns the ability to understand how a model produces
predictions.

Highly interpretable models may include:

\[
\boxed{
\text{small linear models},
}
\]

and

\[
\boxed{
\text{small decision trees}.
}
\]

More complex ensembles can be harder to summarize directly.

%%%%%%%%%%%%%%%%%%%%%%%%%%%%%%%%%%%%%%%%%%%%%%%%%%%%%%%%%%%%
\subsection{Partial Dependence Preview}
\label{subsec:partial-dependence-preview}
%%%%%%%%%%%%%%%%%%%%%%%%%%%%%%%%%%%%%%%%%%%%%%%%%%%%%%%%%%%%

Partial dependence summarizes average model predictions as one or more
features vary.

For feature \(X_j\),

\[
\boxed{
PD_j(x)
=
\mathbb{E}_{\mathbf{X}_{-j}}
[
\widehat{f}
(
x,\mathbf{X}_{-j}
)
].
}
\]

The expectation is approximated using observed predictor values.

Strongly correlated predictors can complicate interpretation.

%%%%%%%%%%%%%%%%%%%%%%%%%%%%%%%%%%%%%%%%%%%%%%%%%%%%%%%%%%%%
\subsection{Individual Conditional Expectation Preview}
\label{subsec:ice-preview}
%%%%%%%%%%%%%%%%%%%%%%%%%%%%%%%%%%%%%%%%%%%%%%%%%%%%%%%%%%%%

Individual conditional expectation, or ICE, plots show how a model's
prediction changes with one predictor for each individual observation.

Unlike partial dependence, ICE displays heterogeneous prediction
responses rather than only their average.

%%%%%%%%%%%%%%%%%%%%%%%%%%%%%%%%%%%%%%%%%%%%%%%%%%%%%%%%%%%%
\subsection{Feature Importance Is Not Causality}
\label{subsec:importance-not-causality}
%%%%%%%%%%%%%%%%%%%%%%%%%%%%%%%%%%%%%%%%%%%%%%%%%%%%%%%%%%%%

A variable may be highly predictive because it:

\[
\boxed{
\text{causes the response},
}
\]

\[
\boxed{
\text{is caused by the response},
}
\]

\[
\boxed{
\text{shares a common cause},
}
\]

or

\[
\boxed{
\text{acts as a proxy for another variable}.
}
\]

Therefore predictive importance does not establish causal importance.

%%%%%%%%%%%%%%%%%%%%%%%%%%%%%%%%%%%%%%%%%%%%%%%%%%%%%%%%%%%%
\subsection{Distribution Shift}
\label{subsec:distribution-shift}
%%%%%%%%%%%%%%%%%%%%%%%%%%%%%%%%%%%%%%%%%%%%%%%%%%%%%%%%%%%%

A model is trained under one data distribution but may later be applied
under another.

This is distribution shift.

Examples include:

\[
\boxed{
\text{changing populations},
}
\]

\[
\boxed{
\text{new geographic regions},
}
\]

\[
\boxed{
\text{changing technology},
}
\]

or

\[
\boxed{
\text{changing economic conditions}.
}
\]

Performance can deteriorate even when the original model was well
validated.

%%%%%%%%%%%%%%%%%%%%%%%%%%%%%%%%%%%%%%%%%%%%%%%%%%%%%%%%%%%%
\subsection{Covariate Shift}
\label{subsec:covariate-shift}
%%%%%%%%%%%%%%%%%%%%%%%%%%%%%%%%%%%%%%%%%%%%%%%%%%%%%%%%%%%%

Covariate shift occurs when

\[
\boxed{
P_{train}(\mathbf{X})
\neq
P_{test}(\mathbf{X})
}
\]

while the conditional relation

\[
P(Y\mid\mathbf{X})
\]

is approximately unchanged.

Reweighting or domain-adaptation methods may be useful.

%%%%%%%%%%%%%%%%%%%%%%%%%%%%%%%%%%%%%%%%%%%%%%%%%%%%%%%%%%%%
\subsection{Concept Drift}
\label{subsec:concept-drift}
%%%%%%%%%%%%%%%%%%%%%%%%%%%%%%%%%%%%%%%%%%%%%%%%%%%%%%%%%%%%

Concept drift occurs when the conditional relationship itself changes:

\[
\boxed{
P_{train}(Y\mid\mathbf{X})
\neq
P_{future}(Y\mid\mathbf{X}).
}
\]

This is especially important in systems operating over time.

Models may require monitoring and updating.

%%%%%%%%%%%%%%%%%%%%%%%%%%%%%%%%%%%%%%%%%%%%%%%%%%%%%%%%%%%%
\subsection{Fairness Preview}
\label{subsec:fairness-preview}
%%%%%%%%%%%%%%%%%%%%%%%%%%%%%%%%%%%%%%%%%%%%%%%%%%%%%%%%%%%%

Prediction quality can differ across groups.

Possible fairness criteria involve:

\[
\boxed{
\text{error rates},
}
\]

\[
\boxed{
\text{calibration},
}
\]

\[
\boxed{
\text{selection rates},
}
\]

or

\[
\boxed{
\text{decision consequences}.
}
\]

Different fairness definitions may conflict mathematically.

The appropriate criterion depends on the application and ethical
context.

%%%%%%%%%%%%%%%%%%%%%%%%%%%%%%%%%%%%%%%%%%%%%%%%%%%%%%%%%%%%
\subsection{Statistical Learning Workflow}
\label{subsec:statistical-learning-workflow}
%%%%%%%%%%%%%%%%%%%%%%%%%%%%%%%%%%%%%%%%%%%%%%%%%%%%%%%%%%%%

A practical workflow is:

\begin{enumerate}
    \item define the prediction or discovery objective;
    \item identify the observational unit;
    \item identify the response and predictors;
    \item create an appropriate train-validation-test strategy;
    \item inspect missing values and measurement quality;
    \item preprocess predictors using training data only;
    \item establish a simple baseline model;
    \item choose candidate model families;
    \item define evaluation metrics before tuning;
    \item tune hyperparameters through cross-validation;
    \item compare models using held-out performance;
    \item inspect calibration and error structure;
    \item examine interpretability and variable importance when useful;
    \item check performance across relevant subgroups;
    \item assess sensitivity to distribution shift;
    \item evaluate once on the untouched test set;
    \item document the entire modeling pipeline.
\end{enumerate}

%%%%%%%%%%%%%%%%%%%%%%%%%%%%%%%%%%%%%%%%%%%%%%%%%%%%%%%%%%%%
\subsection{Common Errors}
\label{subsec:statistical-learning-common-errors}
%%%%%%%%%%%%%%%%%%%%%%%%%%%%%%%%%%%%%%%%%%%%%%%%%%%%%%%%%%%%

\subsubsection{Evaluating Only Training Error}

Training performance measures how well the model fits observations it has
already seen.

It does not estimate future generalization reliably.

%%%%%%%%%%%%%%%%%%%%%%%%%%%%%%%%%%%%%%%%%%%%%%%%%%%%%%%%%%%%
\subsubsection{Tuning on the Test Set}

Repeatedly choosing models according to test performance converts the
test set into validation data.

The reported test performance then becomes optimistically biased.

%%%%%%%%%%%%%%%%%%%%%%%%%%%%%%%%%%%%%%%%%%%%%%%%%%%%%%%%%%%%
\subsubsection{Preprocessing Before Splitting}

Scaling, imputation, feature selection, or dimension reduction using all
data leaks information from validation or test observations.

%%%%%%%%%%%%%%%%%%%%%%%%%%%%%%%%%%%%%%%%%%%%%%%%%%%%%%%%%%%%
\subsubsection{Selecting Features Before Cross-Validation}

Feature selection must occur inside each training fold.

Otherwise the held-out fold indirectly influences the selected predictor
set.

%%%%%%%%%%%%%%%%%%%%%%%%%%%%%%%%%%%%%%%%%%%%%%%%%%%%%%%%%%%%
\subsubsection{Assuming More Complex Means Better}

Greater flexibility reduces training error but may increase test error.

Model complexity must be justified by out-of-sample performance.

%%%%%%%%%%%%%%%%%%%%%%%%%%%%%%%%%%%%%%%%%%%%%%%%%%%%%%%%%%%%
\subsubsection{Assuming More Predictors Always Improve Prediction}

Irrelevant predictors can increase variance, computation, and
overfitting.

High-dimensional models often require regularization.

%%%%%%%%%%%%%%%%%%%%%%%%%%%%%%%%%%%%%%%%%%%%%%%%%%%%%%%%%%%%
\subsubsection{Ignoring Predictor Scale}

Distance-based methods and coefficient penalties depend on scale.

Unstandardized predictors can dominate the model arbitrarily.

%%%%%%%%%%%%%%%%%%%%%%%%%%%%%%%%%%%%%%%%%%%%%%%%%%%%%%%%%%%%
\subsubsection{Interpreting Lasso Selection as Certain Scientific Relevance}

Lasso-selected variables can change across samples, especially under
strong predictor correlation.

Selection does not establish causality or stable scientific importance.

%%%%%%%%%%%%%%%%%%%%%%%%%%%%%%%%%%%%%%%%%%%%%%%%%%%%%%%%%%%%
\subsubsection{Using Accuracy for Highly Imbalanced Data}

High accuracy can occur even when the minority class is never correctly
identified.

Use metrics matched to the decision problem.

%%%%%%%%%%%%%%%%%%%%%%%%%%%%%%%%%%%%%%%%%%%%%%%%%%%%%%%%%%%%
\subsubsection{Interpreting AUC as Calibration}

AUC measures ranking discrimination.

It does not establish that predicted probabilities are numerically
accurate.

%%%%%%%%%%%%%%%%%%%%%%%%%%%%%%%%%%%%%%%%%%%%%%%%%%%%%%%%%%%%
\subsubsection{Choosing a Probability Threshold Automatically at \(0.5\)}

The optimal decision threshold depends on prevalence and consequences of
errors.

%%%%%%%%%%%%%%%%%%%%%%%%%%%%%%%%%%%%%%%%%%%%%%%%%%%%%%%%%%%%
\subsubsection{Interpreting Feature Importance as Causal Effect}

Importance reflects predictive dependence under the fitted algorithm.

It does not identify intervention effects.

%%%%%%%%%%%%%%%%%%%%%%%%%%%%%%%%%%%%%%%%%%%%%%%%%%%%%%%%%%%%
\subsubsection{Using Random Cross-Validation for Dependent Data}

Time series, clustered data, repeated measures, and spatial data require
resampling that respects the dependence structure.

%%%%%%%%%%%%%%%%%%%%%%%%%%%%%%%%%%%%%%%%%%%%%%%%%%%%%%%%%%%%
\subsubsection{Ignoring Distribution Shift}

A model validated on one population may fail in another.

Generalization claims must match the deployment population.

%%%%%%%%%%%%%%%%%%%%%%%%%%%%%%%%%%%%%%%%%%%%%%%%%%%%%%%%%%%%
\subsubsection{Reporting Only the Best Cross-Validation Score}

Model comparison should account for variability across folds or repeated
resamples.

Tiny performance differences may not be practically meaningful.

%%%%%%%%%%%%%%%%%%%%%%%%%%%%%%%%%%%%%%%%%%%%%%%%%%%%%%%%%%%%
\subsubsection{Assuming Unsupervised Clusters Are True Natural Groups}

Clustering always produces structure according to the chosen algorithm,
distance, scaling, and number of clusters.

The resulting groups require external validation and scientific
interpretation.

%%%%%%%%%%%%%%%%%%%%%%%%%%%%%%%%%%%%%%%%%%%%%%%%%%%%%%%%%%%%
\subsubsection{Ignoring Random Initialization}

Algorithms such as \(k\)-means can converge to different local optima.

Multiple starts should be considered.

%%%%%%%%%%%%%%%%%%%%%%%%%%%%%%%%%%%%%%%%%%%%%%%%%%%%%%%%%%%%
\subsection{Applications}
\label{subsec:statistical-learning-applications}
%%%%%%%%%%%%%%%%%%%%%%%%%%%%%%%%%%%%%%%%%%%%%%%%%%%%%%%%%%%%

\begin{applicationbox}

\textbf{Environmental Science.}
Statistical learning can predict environmental measurements, classify
ecosystem conditions, identify spatial patterns, and model complex sensor
data.

\medskip

\textbf{Medicine.}
Classification and survival-learning models can predict diagnosis,
treatment response, recurrence, and clinical risk.

\medskip

\textbf{Business.}
Models predict customer churn, demand, sales, fraud, and purchasing
behavior.

\medskip

\textbf{Finance.}
Regularization, trees, ensembles, and classification methods are used for
credit risk, return prediction, and fraud detection.

\medskip

\textbf{Engineering.}
Statistical learning predicts component failure, identifies anomalies,
models sensor signals, and supports predictive maintenance.

\medskip

\textbf{Education.}
Models can predict student outcomes, identify risk patterns, and analyze
program engagement.

\medskip

\textbf{Biology.}
High-dimensional methods analyze genomic, ecological, and biomedical
measurements.

\medskip

\textbf{Operations.}
Predictive models support forecasting, scheduling, maintenance, routing,
inventory, and resource allocation.

\end{applicationbox}

%%%%%%%%%%%%%%%%%%%%%%%%%%%%%%%%%%%%%%%%%%%%%%%%%%%%%%%%%%%%
\subsection{Exercises}
\label{subsec:statistical-learning-exercises}
%%%%%%%%%%%%%%%%%%%%%%%%%%%%%%%%%%%%%%%%%%%%%%%%%%%%%%%%%%%%

\begin{exercise}[1]
Define supervised learning.
\end{exercise}

\begin{exercise}[1]
Define unsupervised learning.
\end{exercise}

\begin{exercise}[1]
Distinguish regression and classification.
\end{exercise}

\begin{exercise}[1]
Define training error.
\end{exercise}

\begin{exercise}[1]
Define generalization error.
\end{exercise}

\begin{exercise}[1]
Define overfitting.
\end{exercise}

\begin{exercise}[1]
Define regularization.
\end{exercise}

\begin{exercise}[1]
Define cross-validation.
\end{exercise}

\begin{exercise}[1]
Define data leakage.
\end{exercise}

\begin{exercise}[1]
Define ridge regression.
\end{exercise}

\begin{exercise}[1]
Define lasso regression.
\end{exercise}

\begin{exercise}[1]
Define a decision tree.
\end{exercise}

\begin{exercise}[1]
Define bagging.
\end{exercise}

\begin{exercise}[1]
Define a random forest.
\end{exercise}

\begin{exercise}[1]
Define boosting.
\end{exercise}

\begin{exercise}[1]
Define a support vector.
\end{exercise}

\begin{exercise}[1]
Define \(k\)-means clustering.
\end{exercise}

\begin{exercise}[2]
Suppose a model has training MSE \(2\) and test MSE \(15\). What problem
might this indicate?
\end{exercise}

\begin{exercise}[2]
Suppose both training and test errors are very high. What problem might
this indicate?
\end{exercise}

\begin{exercise}[2]
For centered one-variable ridge regression with

\[
X^TX=8,
\qquad
X^Ty=12,
\qquad
\lambda=4,
\]

compute the ridge coefficient.
\end{exercise}

\begin{exercise}[2]
Explain what happens to ridge coefficients as

\[
\lambda\to\infty.
\]
\end{exercise}

\begin{exercise}[2]
Explain what happens to ridge regression as

\[
\lambda\to0.
\]
\end{exercise}

\begin{exercise}[2]
State one major difference between ridge and lasso regression.
\end{exercise}

\begin{exercise}[2]
Suppose a binary classifier has

\[
TP=80,
\quad
FN=20.
\]

Compute sensitivity.
\end{exercise}

\begin{exercise}[2]
Suppose

\[
TN=90,
\quad
FP=10.
\]

Compute specificity.
\end{exercise}

\begin{exercise}[2]
Suppose

\[
TP=40,
\quad
FP=10.
\]

Compute precision.
\end{exercise}

\begin{exercise}[2]
Compute the \(F_1\) score when precision and recall are both \(0.8\).
\end{exercise}

\begin{exercise}[2]
A \(k\)-means problem has points

\[
1,\ 2,\ 10,\ 11
\]

in one dimension and \(K=2\). Identify an obvious two-cluster solution
and its cluster means.
\end{exercise}

\begin{exercise}[2]
Explain why KNN should usually use standardized predictors when units
differ substantially.
\end{exercise}

\begin{exercise}[3]
Derive the squared-error bias--variance decomposition.
\end{exercise}

\begin{exercise}[3]
Derive the ridge-regression estimator

\[
\widehat{\boldsymbol{\beta}}
=
(
X^TX+\lambda I
)^{-1}
X^T\mathbf{y}.
\]
\end{exercise}

\begin{exercise}[3]
Explain why ridge regression is useful under multicollinearity.
\end{exercise}

\begin{exercise}[3]
Explain geometrically why lasso can produce exact zeros.
\end{exercise}

\begin{exercise}[3]
Compare ridge, lasso, and elastic net.
\end{exercise}

\begin{exercise}[3]
Explain why feature selection must occur inside cross-validation folds.
\end{exercise}

\begin{exercise}[3]
Compare \(k\)-fold cross-validation and LOOCV.
\end{exercise}

\begin{exercise}[3]
Explain why nested cross-validation may be needed for unbiased model
assessment.
\end{exercise}

\begin{exercise}[3]
Explain why AUC and calibration measure different properties.
\end{exercise}

\begin{exercise}[3]
Explain why a threshold of \(0.5\) may be inappropriate in an imbalanced
classification problem.
\end{exercise}

\begin{exercise}[3]
Derive the majority-vote rule for KNN classification.
\end{exercise}

\begin{exercise}[3]
Explain the bias--variance effect of increasing \(k\) in KNN.
\end{exercise}

\begin{exercise}[3]
Explain how the curse of dimensionality affects nearest-neighbor methods.
\end{exercise}

\begin{exercise}[3]
Derive the optimal terminal-node prediction in a regression tree under
squared-error loss.
\end{exercise}

\begin{exercise}[3]
Compare Gini impurity and entropy.
\end{exercise}

\begin{exercise}[3]
Explain why single trees have high variance.
\end{exercise}

\begin{exercise}[3]
Explain how bagging reduces tree variance.
\end{exercise}

\begin{exercise}[3]
Explain why random forests decorrelate trees.
\end{exercise}

\begin{exercise}[3]
Derive the variance formula for an average of correlated predictors.
\end{exercise}

\begin{exercise}[3]
Explain how out-of-bag error is constructed.
\end{exercise}

\begin{exercise}[3]
Explain how gradient boosting fits residuals under squared-error loss.
\end{exercise}

\begin{exercise}[3]
Compare bagging and boosting.
\end{exercise}

\begin{exercise}[3]
Explain the margin concept in support-vector classification.
\end{exercise}

\begin{exercise}[3]
Explain the role of the kernel trick.
\end{exercise}

\begin{exercise}[3]
Show that a linear kernel reproduces the ordinary inner product.
\end{exercise}

\begin{exercise}[3]
Compare PCA and principal-components regression.
\end{exercise}

\begin{exercise}[3]
Explain why PCR can discard a low-variance but highly predictive
direction.
\end{exercise}

\begin{exercise}[3]
Explain how PLS differs conceptually from PCR.
\end{exercise}

\begin{exercise}[3]
Derive the \(k\)-means update for a cluster centroid.
\end{exercise}

\begin{exercise}[3]
Compare single, complete, and average linkage.
\end{exercise}

\begin{exercise}[3]
Explain why different cluster scaling choices can produce different
results.
\end{exercise}

\begin{exercise}[3]
Explain why variable importance does not imply causality.
\end{exercise}

\begin{exercise}[3]
Explain the difference between covariate shift and concept drift.
\end{exercise}

\begin{exercise}[4]
Study theoretical generalization bounds for empirical risk minimization.
\end{exercise}

\begin{exercise}[4]
Study VC dimension and classification complexity.
\end{exercise}

\begin{exercise}[4]
Investigate structural risk minimization.
\end{exercise}

\begin{exercise}[4]
Derive ridge regression through singular-value decomposition.
\end{exercise}

\begin{exercise}[4]
Study lasso optimality using subgradient conditions.
\end{exercise}

\begin{exercise}[4]
Derive coordinate descent for lasso regression.
\end{exercise}

\begin{exercise}[4]
Study elastic-net solution paths.
\end{exercise}

\begin{exercise}[4]
Investigate post-selection inference.
\end{exercise}

\begin{exercise}[4]
Study stability selection.
\end{exercise}

\begin{exercise}[4]
Derive analytic LOOCV formulas for ordinary least squares.
\end{exercise}

\begin{exercise}[4]
Study bootstrap estimates of prediction error.
\end{exercise}

\begin{exercise}[4]
Investigate random-forest consistency.
\end{exercise}

\begin{exercise}[4]
Study boosting as functional gradient descent.
\end{exercise}

\begin{exercise}[4]
Derive the dual optimization problem for support-vector machines.
\end{exercise}

\begin{exercise}[4]
Study Mercer's theorem for kernel methods.
\end{exercise}

\begin{exercise}[4]
Investigate kernel ridge regression.
\end{exercise}

\begin{exercise}[4]
Study reproducing-kernel Hilbert spaces.
\end{exercise}

\begin{exercise}[4]
Investigate nonlinear dimension-reduction methods.
\end{exercise}

\begin{exercise}[4]
Study spectral clustering.
\end{exercise}

\begin{exercise}[4]
Investigate Gaussian-mixture clustering.
\end{exercise}

\begin{exercise}[4]
Study model-based clustering.
\end{exercise}

\begin{exercise}[4]
Investigate calibration methods for probabilistic classifiers.
\end{exercise}

\begin{exercise}[4]
Study conformal prediction.
\end{exercise}

\begin{exercise}[4]
Investigate learning under distribution shift.
\end{exercise}

\begin{exercise}[4]
Study mathematical tradeoffs among competing fairness criteria.
\end{exercise}

%%%%%%%%%%%%%%%%%%%%%%%%%%%%%%%%%%%%%%%%%%%%%%%%%%%%%%%%%%%%
\subsection{Conceptual Summary}
\label{subsec:statistical-learning-summary}
%%%%%%%%%%%%%%%%%%%%%%%%%%%%%%%%%%%%%%%%%%%%%%%%%%%%%%%%%%%%

Statistical learning seeks a function

\[
\boxed{
\widehat{f}
}
\]

that performs well on future observations.

The population risk is

\[
\boxed{
R(f)
=
\mathbb{E}
[
L
(
Y,f(\mathbf{X})
)
].
}
\]

The empirical risk is

\[
\boxed{
\widehat{R}_n(f)
=
\frac1n
\sum_{i=1}^{n}
L
(
Y_i,f(\mathbf{X}_i)
).
}
\]

Generalization requires controlling the tradeoff between

\[
\boxed{
\text{bias}
}
\]

and

\[
\boxed{
\text{variance}.
}
\]

Regularization modifies the fitting objective through

\[
\boxed{
\text{Loss}
+
\lambda
\cdot
\text{Penalty}.
}
\]

Ridge regression uses

\[
\boxed{
\sum_j\beta_j^2,
}
\]

while lasso uses

\[
\boxed{
\sum_j|\beta_j|.
}
\]

Cross-validation estimates out-of-sample performance and selects tuning
parameters.

Tree-based methods partition predictor space.

Bagging and random forests average many trees.

Boosting builds learners sequentially.

Support-vector machines use margins and kernels.

Unsupervised methods such as PCA and clustering discover structure
without response labels.

Statistical learning therefore combines

\[
\boxed{
\text{statistics}
+
\text{optimization}
+
\text{prediction}
+
\text{regularization}
+
\text{computation}.
}
\]

%%%%%%%%%%%%%%%%%%%%%%%%%%%%%%%%%%%%%%%%%%%%%%%%%%%%%%%%%%%%
\subsection{Section Connections}
\label{subsec:statistical-learning-connections}
%%%%%%%%%%%%%%%%%%%%%%%%%%%%%%%%%%%%%%%%%%%%%%%%%%%%%%%%%%%%

\chapterconnection{
Statistical Learning extends regression, classification, multivariate
statistics, and decision theory toward predictive modeling and
high-dimensional data. Ridge, lasso, and elastic net extend the
regularization ideas introduced in Section~\ref{sec:regression};
principal-components regression connects directly with
Section~\ref{sec:multivariate-statistics}; and cross-validation provides
a general framework for tuning predictive models. Trees, random forests,
boosting, support-vector machines, clustering, and kernel methods broaden
the class of available models beyond traditional parametric forms. The
next section, Computational Statistics, develops the numerical,
simulation, resampling, optimization, and Monte Carlo methods required to
implement modern statistical inference and learning at scale.
}

%%%%%%%%%%%%%%%%%%%%%%%%%%%%%%%%%%%%%%%%%%%%%%%%%%%%%%%%%%%%
% SECTION SOURCE TRACKING
%%%%%%%%%%%%%%%%%%%%%%%%%%%%%%%%%%%%%%%%%%%%%%%%%%%%%%%%%%%%

%------------------------------------------------------------
% 8.18 Statistical Learning
%------------------------------------------------------------
%
% Source basis:
%   - Standard statistical-learning theory.
%   - Standard predictive modeling and regularization.
%   - Standard classification, tree, ensemble, kernel, and clustering
%     methods.
%
% Used for:
%   - supervised and unsupervised learning;
%   - regression and classification;
%   - training, validation, and test data;
%   - prediction and generalization;
%   - loss and risk;
%   - empirical risk minimization;
%   - irreducible error;
%   - bias--variance tradeoff;
%   - underfitting and overfitting;
%   - model complexity;
%   - regularization;
%   - ridge regression;
%   - lasso regression;
%   - elastic net;
%   - feature selection;
%   - cross-validation;
%   - LOOCV;
%   - nested cross-validation;
%   - data leakage;
%   - classification metrics;
%   - ROC and AUC;
%   - calibration;
%   - logistic classification;
%   - k-nearest neighbors;
%   - curse of dimensionality;
%   - decision trees;
%   - regression and classification trees;
%   - Gini impurity;
%   - entropy;
%   - pruning;
%   - bagging;
%   - random forests;
%   - out-of-bag estimation;
%   - variable importance;
%   - boosting;
%   - gradient boosting;
%   - support-vector classification;
%   - kernel methods;
%   - principal-components regression;
%   - partial least squares preview;
%   - high-dimensional learning;
%   - sparsity;
%   - k-means clustering;
%   - hierarchical clustering;
%   - linkage criteria;
%   - distance measures;
%   - ensembles;
%   - stacking preview;
%   - hyperparameter tuning;
%   - learning curves;
%   - imbalanced classification;
%   - interpretability;
%   - partial dependence and ICE previews;
%   - distribution shift;
%   - covariate shift;
%   - concept drift;
%   - fairness preview;
%   - applications and exercises.
%
% Notes:
%   - Computational Statistics is developed next in Section 8.19.
%   - Optimization theory is developed more deeply in Chapter 10.
%   - Numerical and computational implementation is developed more deeply
%     in Chapter 11.
%
%%%%%%%%%%%%%%%%%%%%%%%%%%%%%%%%%%%%%%%%%%%%%%%%%%%%%%%%%%%%